% Guia do Professor — Meteorologia Prática
% Baseado em R. Stull, Practical Meteorology (CC BY-NC-SA 4.0)
\documentclass[11pt]{article}
\usepackage{../comum/preambulo}

\title{Guia do Professor\\ {\large Meteorologia Prática ---
curso baseado em R.~Stull, \emph{Practical Meteorology}}}
\author{}
\date{Versão de trabalho}

\begin{document}
\maketitle
\thispagestyle{fancy}

\section*{Sobre este guia}

Este guia acompanha, capítulo a capítulo, o livro-texto
\emph{Practical Meteorology: An Algebra-based Survey of Atmospheric Science}
(R.~Stull, UBC, v1.02b, CC BY-NC-SA 4.0). Para cada capítulo há:
objetivos de aprendizagem, pré-requisitos, ritmo sugerido, armadilhas
conceituais comuns, sugestões de problemas do livro, e o roteiro de uso do
notebook de simulação em aula.

\paragraph{Filosofia do curso.} O livro é deliberadamente algébrico; nosso
público é de estudantes de física com cálculo. As notas de aula mostram
\emph{as duas formas} de cada lei — a algébrica do livro (para leitura e
provas com consulta ao livro) e a diferencial (para conexão com mecânica dos
fluidos e termodinâmica). Os notebooks fecham o ciclo: cada equação central
é verificada numericamente contra dados ou contra a atmosfera padrão.

\paragraph{Estrutura de cada capítulo do curso.}
\begin{itemize}
  \item \texttt{notas/capNN\_notas.tex} — roteiro de quadro-negro (o que
        escrever, em que ordem, com deduções completas);
  \item \texttt{slides/capNN\_slides.tex} — apoio visual (figuras do livro,
        tabelas, sínteses); o quadro carrega as deduções, os slides carregam
        as imagens;
  \item \texttt{notebooks/capNN\_*.ipynb} — 2 a 4 casos numéricos por
        capítulo, para demonstração em aula e exploração pelos alunos;
  \item este guia — a ``regência'' de tudo.
\end{itemize}

\paragraph{Ambiente computacional.} \texttt{conda env create -f
environment.yml} cria o ambiente \texttt{meteo} com MetPy, climlab,
ambiance, pyrcel, tcpyPI, Py-ART e Satpy. Os notebooks declaram no topo os
pacotes que usam.

\bigskip
\hrule
\bigskip

% Guia do professor — Capítulo 1: Fundamentos da Atmosfera
\section{Capítulo 1 --- Fundamentos da Atmosfera}

\subsection*{Visão geral e ritmo}

Capítulo de fundação: convenções, estado termodinâmico $(P,T,\rho)$,
estrutura vertical, lei dos gases, hidrostática e equação hipsométrica.
\textbf{Ritmo sugerido: 2 aulas de 2\,h.}

\begin{itemize}
  \item \textbf{Aula 1} — convenções meteorológicas, referenciais,
        $(P,T,\rho)$, estrutura vertical e atmosfera padrão.
        Fechar com o Caso~1 do notebook (10\,min de projeção).
  \item \textbf{Aula 2} — lei dos gases + temperatura virtual, dedução da
        hidrostática, dedução da hipsométrica, processos.
        Fechar com Casos~2--4 (15\,min).
\end{itemize}

\subsection*{Objetivos de aprendizagem}

Ao final, o aluno deve ser capaz de:
\begin{enumerate}
  \item converter unidades de $T$ e $P$ e usar a convenção de direção do vento;
  \item aplicar $P=\rho \Re_d T_v$, explicando fisicamente por que ar úmido é
        menos denso e o papel de $T_v$;
  \item deduzir $dP/dz=-\rho g$ por balanço de forças e estimar variações de
        pressão com a altura;
  \item deduzir e aplicar a equação hipsométrica; interpretar espessura como
        temperatura média da camada;
  \item descrever as camadas da atmosfera e o porquê físico do perfil $T(z)$
        de cada uma;
  \item rodar o notebook do capítulo e interpretar os gráficos gerados.
\end{enumerate}

\subsection*{Pré-requisitos a reativar}

Gás ideal na forma molecular ($PV=Nk_BT$), integração de $1/x$, e noção de
massa molar média. Alunos de física raramente viram a forma
\emph{por unidade de massa} $P=\rho\Re T$ — vale 5 minutos de ponte explícita
(está na caixa ``forma diferencial'' das notas).

\subsection*{Armadilhas conceituais frequentes}

\begin{itemize}
  \item \textbf{Direção do vento invertida}: alunos assumem que ``vento de
        oeste'' vai \emph{para} oeste. Insistir: nomeia-se a origem.
  \item \textbf{kPa vs.\ hPa}: o livro usa kPa, a operação usa hPa; erros de
        fator 10 são a falha mais comum nas listas. Treinar as duas.
  \item \textbf{$T_v$ ``esfriando''}: $T_v \ge T$ para ar sem condensados;
        parte dos alunos memoriza ao contrário. O Caso~4 do notebook mostra a
        magnitude real ($\lesssim 3$--4\,K nos trópicos).
  \item \textbf{Hidrostática como lei exata}: enfatizar que é aproximação de
        equilíbrio — dentro de uma corrente ascendente de tempestade ela
        falha (gancho para o Cap.~14).
  \item \textbf{Escalas de altura}: $H_p \ne H_\rho$ confunde; a pergunta
        ``por quê?'' rende boa discussão (resposta: $T$ varia com $z$).
\end{itemize}

\subsection*{Demonstração de baixo custo}

Barômetro do celular (apps gratuitos leem o sensor MEMS): medir $P$ no térreo
e no último andar do prédio; comparar com $\Delta P=-\rho g \Delta z$.
Rende 10 minutos memoráveis e conecta imediatamente com a Aula 2.

\subsection*{Problemas sugeridos do livro (seção 1.12)}

\begin{itemize}
  \item \textbf{Aquecimento (Apply)}: A1, A5 (conversões), A9 (direção do
        vento), A14--A16 (gás ideal e $T_v$), A19 (hidrostática),
        A22 (hipsométrica).
  \item \textbf{Aprofundamento (Evaluate \& Analyze)}: E5 (escala de altura),
        E12 (espessura vs.\ temperatura) — bons para entrega escrita.
  \item \textbf{Síntese (Synthesize)}: S1 (``e se a atmosfera fosse
        isotérmica?'') — casa exatamente com a dedução em caixa vermelha das
        notas e com o Caso~1 do notebook.
\end{itemize}
\emph{Obs.: confira a numeração na sua cópia do livro (v1.02b) antes de
publicar a lista.}

\subsection*{Uso do notebook \texttt{cap01\_fundamentos.ipynb}}

\begin{center}
\begin{tabular}{lll}
\hline
Caso & Conteúdo & Momento sugerido \\
\hline
1 & Atmosfera padrão ICAO anotada (\texttt{ambiance}); teste log &
fim da Aula 1 \\
2 & Integração numérica da EDO hidrostática (\texttt{solve\_ivp}) &
Aula 2, após hidrostática \\
3 & Altimetria barométrica; espessura 1000--500 vs.\ $\bar T_v$ &
Aula 2, após hipsométrica \\
4 & Sondagem real (\texttt{MetPy}, Skew-T) vs.\ atmosfera padrão &
Aula 2 ou tarefa de casa \\
5 & Temperatura virtual e densidade do ar úmido &
Aula 2, após $T_v$ \\
\hline
\end{tabular}
\end{center}

O Caso~2 merece destaque pedagógico: é a primeira ``previsão numérica em
miniatura'' do curso (lei diferencial + condição inicial + integrador) e
planta a semente do Cap.~20.

Sugestão de tarefa: pedir que cada aluno rode o Caso~4 com a sondagem mais
recente de uma estação brasileira (ex.: SBMT/Campo de Marte via Wyoming) e
comente três diferenças em relação à atmosfera padrão.

\subsection*{Exercícios complementares e soluções}

As notas de aula trazem 5 exercícios teóricos (T1--T5, para entrega escrita)
e 4 numéricos (N1--N4, com células reservadas no fim de
\texttt{cap01\_fundamentos.ipynb}). O notebook
\texttt{cap01\_solucoes.ipynb} --- \emph{não distribuir antes da entrega} ---
resolve todos: os teóricos com verificação simbólica em \texttt{sympy}
(EDOs, expansões em série, limites) e os numéricos com código completo e
gráficos. T2$\leftrightarrow$N1 formam um par teoria--numérica proposital:
o aluno deduz o perfil politrópico e depois o confere com o integrador.

\subsection*{Conexões com o resto do curso}

Gases ideais e $T_v$ $\to$ estabilidade (Cap.~5); hipsométrica $\to$ cartas
de espessura (Cap.~9) e vento térmico (Cap.~11); camada limite $\to$
Cap.~18; falha da hidrostática $\to$ tempestades (Cap.~14).

\clearpage
% Guia do professor — Capítulo 2: Radiação Solar e Infravermelha
\section{Capítulo 2 --- Radiação Solar e Infravermelha}

\subsection*{Visão geral e ritmo}

O capítulo energético do curso: geometria Sol--Terra, leis de corpo negro,
lei de Beer e balanço radiativo, culminando na temperatura efetiva e no
efeito estufa. \textbf{Ritmo sugerido: 2 aulas de 2\,h.}

\begin{itemize}
  \item \textbf{Aula 1} — motivação (radiação como combustível), declinação,
        elevação solar, duração do dia, fluxo e inverso do quadrado.
        Fechar com o Caso~2 (mapa de insolação — 10\,min).
  \item \textbf{Aula 2} — Planck, Wien, Stefan--Boltzmann, Kirchhoff, Beer,
        balanço na superfície, $T_e = 255$\,K e o efeito estufa.
        Fechar com o Caso~4 (climlab — 15\,min).
\end{itemize}

\subsection*{Objetivos de aprendizagem}

\begin{enumerate}
  \item calcular declinação, elevação solar e duração do dia para qualquer
        latitude e data;
  \item aplicar Planck/Wien/Stefan--Boltzmann e justificar a separação
        onda curta / onda longa;
  \item deduzir e aplicar a lei de Beer (forma diferencial e integrada);
  \item calcular $T_e$ de um planeta e explicar quantitativamente o efeito
        estufa (modelo de 1 e $N$ camadas);
  \item montar o balanço radiativo de superfície com albedo e IV atmosférico.
\end{enumerate}

\subsection*{Pré-requisitos a reativar}

Trigonometria esférica mínima (a eq.\ da elevação pode ser \emph{dada});
integral de $x^3/(e^x-1)$ apenas citada (solução simbólica no notebook de
soluções); lei dos gases não é usada — capítulo independente do Cap.~1
exceto pela estrutura vertical.

\subsection*{Armadilhas conceituais frequentes}

\begin{itemize}
  \item \textbf{``Verão porque a Terra está mais perto do Sol''}: o periélio
        é em janeiro — e o verão do hemisfério norte é em julho. O que manda
        é $\sin\Psi$ e a duração do dia. Vale perguntar antes de explicar.
  \item \textbf{Fator 4 do $T_e$}: disco intercepta, esfera emite. Alunos
        esquecem e obtêm 279\,K em vez de 255\,K (sem albedo) — apontar que
        os dois erros (fator e albedo) quase se cancelam.
  \item \textbf{$T_e$ vs.\ temperatura da superfície}: $T_e$ é a temperatura
        de \emph{emissão para o espaço} (nível médio $\sim$5\,km), não do
        chão.
  \item \textbf{Wien com $T$ em °C}: erro clássico; reforçar kelvins.
  \item \textbf{Beer para radiação difusa}: a lei vale para o feixe direto;
        céu nublado exige tratamento de espalhamento múltiplo (fora do
        escopo).
\end{itemize}

\subsection*{Demonstração de baixo custo}

Termômetro IV de cozinha (ou câmera térmica de celular): medir o ``céu''
limpo (dá $-20$ a $-50\,°$C — o $I\!\downarrow$ baixo do ar seco) e uma
nuvem (bem mais quente). Conecta emissividade, $I\!\downarrow$ e efeito
estufa em 5 minutos.

\subsection*{Problemas sugeridos do livro (seção 2.9)}

Aquecimento: A2 (declinação), A7 (elevação), A12--A14 (Planck/Wien/S-B),
A19 (Beer), A22 (balanço). Aprofundamento: E8 (insolação diária),
E15 (albedo e $F^*$). Síntese: S3 (``e se o eixo da Terra fosse
$45°$?'') — casa com o Caso~2. \emph{Confira a numeração na sua cópia
(v1.02b) antes de publicar.}

\subsection*{Uso do notebook \texttt{cap02\_radiacao.ipynb}}

\begin{center}
\begin{tabular}{lll}
\hline
Caso & Conteúdo & Momento sugerido \\
\hline
1 & Curvas de Planck Sol × Terra; Wien e S--B verificados & Aula 2, após Planck \\
2 & Elevação solar, duração do dia, mapa de insolação lat × dia & fim da Aula 1 \\
3 & Lei de Beer como EDO; atenuação vs.\ ângulo zenital & Aula 2, após Beer \\
4 & Estufa de $N$ camadas × coluna cinza \texttt{climlab} & fim da Aula 2 \\
\hline
\end{tabular}
\end{center}

O Caso~2 rende o gráfico mais bonito do início do curso (mapa de insolação
latitude × dia); o Caso~4 introduz o \texttt{climlab}, que voltará no
Cap.~21.

\subsection*{Exercícios complementares e soluções}

Nas notas: T1--T5 (Wien e S--B a partir de Planck, $T_e$ planetário, estufa
de 1 camada, massa de ar $1/\sin\Psi$) e N1--N4 (células reservadas no
notebook). Soluções completas em \texttt{cap02\_solucoes.ipynb}
(\emph{não distribuir}): T1--T2 verificados simbolicamente com
\texttt{sympy}, N4 compara a fórmula $(N{+}1)^{1/4}$ com o modelo cinza.
Par teoria--numérica: T4$\leftrightarrow$N4.

\subsection*{Conexões com o resto do curso}

$F^*$ alimenta o balanço de calor da superfície (Cap.~3); insolação
diferencial $\to$ circulação geral (Cap.~11); Beer $\to$ satélites
(Cap.~8); climlab e estufa $\to$ clima (Cap.~21).

\clearpage
% Guia do professor — Capítulo 3: Termodinâmica da Atmosfera
\section{Capítulo 3 --- Termodinâmica da Atmosfera}

\subsection*{Visão geral e ritmo}

O capítulo transforma a 1ª lei da termodinâmica em quatro ferramentas:
gradiente adiabático seco, temperatura potencial, balanço euleriano de
calor e balanço de superfície. \textbf{Ritmo sugerido: 2 aulas de 2\,h.}

\begin{itemize}
  \item \textbf{Aula 1} — calores específicos e latente, 1ª lei na forma
        meteorológica, $\Gamma_d$, temperatura potencial (com o exemplo
        föhn). Fechar com o Caso~1 (10\,min).
  \item \textbf{Aula 2} — derivada material e advecção, balanço euleriano,
        balanço de superfície e Bowen. Fechar com Casos~2--3 (15--20\,min;
        o Caso~3 merece a projeção mais longa).
\end{itemize}

\subsection*{Objetivos de aprendizagem}

\begin{enumerate}
  \item aplicar $\delta q = c_p dT - \alpha dP$ e explicar fisicamente cada
        termo;
  \item deduzir $\Gamma_d = g/c_p$ e distingui-lo do gradiente ambiental;
  \item calcular $\theta$, justificar sua conservação e usá-la para comparar
        parcelas de níveis diferentes;
  \item escrever o balanço euleriano, estimar advecção com números reais e
        entender a derivada material;
  \item repartir $F^*$ em $F_H$, $F_E$, $F_G$ via razão de Bowen e prever o
        caráter do ciclo diurno de uma superfície.
\end{enumerate}

\subsection*{Pré-requisitos a reativar}

1ª lei na forma $dU = \delta Q - PdV$ (física básica); hidrostática e gás
ideal do Cap.~1 (usados nas duas deduções centrais); regra da cadeia em
várias variáveis (derivada material). Alunos de física têm tudo — o novo é
o \emph{vocabulário} meteorológico.

\subsection*{Armadilhas conceituais frequentes}

\begin{itemize}
  \item \textbf{``Esfria porque lá em cima é frio''}: a parcela esfria por
        \emph{expansão}, não por contato. É a confusão número 1 do capítulo.
  \item \textbf{$\Gamma_d$ vs.\ gradiente ambiental}: 9{,}8 é o que a
        parcela faz; 6{,}5 é a média do que a sondagem mostra. Confundi-los
        inviabiliza o Cap.~5.
  \item \textbf{$\theta$ como ``temperatura exótica''}: insistir que é só um
        rótulo conservado — comparar com energia mecânica em queda livre.
  \item \textbf{Sinal da advecção}: vento soprando \emph{do} quente
        \emph{para} o frio dá $-U\,\partial T/\partial x > 0$ (aquece).
        Fazer o desenho do gradiente com setas antes da fórmula.
  \item \textbf{Bowen como propriedade fixa}: $\beta$ depende da umidade
        disponível — o mesmo solo seco ou molhado muda $\beta$ por ordem de
        magnitude.
\end{itemize}

\subsection*{Demonstração de baixo custo}

Borrifador de álcool no termômetro (ou no dorso da mão): queda imediata de
vários graus — calor latente em ação. Alternativa: bomba manual de encher
pneu aquece ao comprimir (adiabática no sentido oposto).

\subsection*{Problemas sugeridos do livro (seção de exercícios do Cap.~3)}

Aquecimento: conversões de calor/temperatura, cálculo de $\theta$ para
vários níveis, advecção com vento e gradiente dados. Aprofundamento:
problemas de balanço de superfície com Bowen; síntese: ``o que aconteceria
se $L_v$ da água fosse 10× menor?''. \emph{Confira a numeração na sua cópia
(v1.02b) antes de publicar a lista.}

\subsection*{Uso do notebook \texttt{cap03\_termodinamica.ipynb}}

\begin{center}
\begin{tabular}{lll}
\hline
Caso & Conteúdo & Momento sugerido \\
\hline
1 & Parcela adiabática (EDO); conservação de $\theta$ & fim da Aula 1 \\
2 & Ciclo diurno: balanço de superfície com Bowen & Aula 2, após Bowen \\
3 & Advecção 1D (EDP, esquema \emph{upwind}) & Aula 2, após advecção \\
4 & $\theta$ da sondagem real; camada de mistura & Aula 2 ou tarefa \\
\hline
\end{tabular}
\end{center}

O Caso~3 é o momento didático mais rico: é a primeira EDP integrada no
curso (``previsão de verdade''), e a difusão numérica visível abre a
conversa sobre por que modelos operacionais usam esquemas sofisticados
(Cap.~20).

\subsection*{Exercícios complementares e soluções}

Nas notas: T1--T5 ($\Gamma_d$ terrestre e marciano, dedução de $\theta$,
$\theta$--entropia, solução geral da advecção, partição de Bowen) e N1--N4
(células reservadas no notebook). Soluções em
\texttt{cap03\_solucoes.ipynb} (\emph{não distribuir}): T2--T4 com
verificação simbólica \texttt{sympy}; N3 quantifica a difusão numérica do
\emph{upwind}. Pares teoria--numérica: T2$\leftrightarrow$N1 e
T4$\leftrightarrow$N3.

\subsection*{Conexões com o resto do curso}

$\Gamma_d$ e $\theta$ $\to$ estabilidade e diagramas termodinâmicos
(Cap.~5); calor latente $\to$ umidade (Cap.~4) e tempestades (Cap.~14);
advecção $\to$ frentes (Cap.~12) e NWP (Cap.~20); balanço de superfície
$\to$ camada limite (Cap.~18).

\clearpage
% Guia do professor — Capítulo 4: Vapor d'Água
\section{Capítulo 4 --- Vapor d'Água}

\subsection*{Visão geral e ritmo}

O dicionário quantitativo da umidade: Clausius--Clapeyron, as sete
variáveis de umidade, LCL, bulbo úmido, água precipitável e o balanço de
água. \textbf{Ritmo sugerido: 2 aulas de 2\,h.}

\begin{itemize}
  \item \textbf{Aula 1} — motivação, dedução e integração de
        Clausius--Clapeyron, a curva $e_s(T)$, o zoológico de variáveis
        (com o exemplo de São Paulo). Fechar com Casos~1--2 (15\,min).
  \item \textbf{Aula 2} — LCL (com dedução da regra dos 125\,m/°C), bulbo
        úmido e limite fisiológico, água precipitável, balanço de água.
        Fechar com Casos~3--4 (10\,min).
\end{itemize}

\subsection*{Objetivos de aprendizagem}

\begin{enumerate}
  \item deduzir e integrar Clausius--Clapeyron; explicar a regra dos
        7\,\%/K e suas consequências climáticas;
  \item converter fluentemente entre $e$, $r$, $q$, $\rho_v$, UR, $T_d$,
        $T_w$ e saber qual usar em cada problema;
  \item explicar por que $r$ é conservada e UR não presta para comparar
        massas de ar;
  \item calcular o LCL pela regra e pelo método exato; ler a base do
        cumulus como higrômetro;
  \item calcular água precipitável de uma sondagem e montar o balanço
        euleriano de água.
\end{enumerate}

\subsection*{Pré-requisitos a reativar}

Gás ideal por espécie (pressões parciais, Cap.~1); $\Gamma_d$ e conservação
de $\theta$ (Cap.~3); potencial de Gibbs é citado mas a dedução de
Clausius--Clapeyron pode ser dada ``por decreto'' com a EDO como ponto de
partida — físicos aceitam bem.

\subsection*{Armadilhas conceituais frequentes}

\begin{itemize}
  \item \textbf{``Ar quente segura mais água''}: linguagem de esponja
        esconde a física — é a \emph{pressão de equilíbrio} $e_s(T)$ que
        cresce; o ar não ``segura'' nada. Vale corrigir explicitamente.
  \item \textbf{UR como medida absoluta}: 90\,\% de UR no inverno pode ser
        menos vapor que 40\,\% no verão. O exemplo numérico cura.
  \item \textbf{$T_d$ vs.\ $T_w$}: alunos trocam; fixar
        $T_d \le T_w \le T$ com o caso saturado degenerando tudo.
  \item \textbf{Fator $\varepsilon$}: esquecem o 0,622 e erram $r$ por
        60\,\%; lembrar a origem molecular (18/29).
  \item \textbf{LCL para ar não aquecido pela superfície}: a regra dos
        125\,m/°C usa $T$ e $T_d$ \emph{da parcela que sobe}; para
        estratiforme o levantamento é outro (Cap.~5).
\end{itemize}

\subsection*{Demonstração de baixo custo}

Psicrômetro artesanal: dois termômetros idênticos, gaze molhada num deles,
abanar 1 minuto. Medir $T$ e $T_w$, calcular $r$ e UR em sala com a equação
psicrométrica (é o exercício N2 ao vivo). Alternativa: lata gelada suando —
a película é o orvalho, a temperatura da lata no início da condensação
$\approx T_d$ do ar da sala.

\subsection*{Problemas sugeridos do livro (seção de exercícios do Cap.~4)}

Aquecimento: cálculo de $e_s$ para várias $T$; conversões $e \to r \to$ UR;
$T_d$ e LCL. Aprofundamento: bulbo úmido e psicrômetro; água precipitável
de perfis dados. Síntese: ``por que o deserto tem a maior amplitude térmica
diária?'' (liga UR, $I\!\downarrow$ do Cap.~2 e Bowen do Cap.~3).
\emph{Confira a numeração na sua cópia (v1.02b).}

\subsection*{Uso do notebook \texttt{cap04\_vapor.ipynb}}

\begin{center}
\begin{tabular}{lll}
\hline
Caso & Conteúdo & Momento sugerido \\
\hline
1 & EDO de Clausius--Clapeyron × fórmula × MetPy; gelo × líquido & Aula 1, após CC \\
2 & Carta psicrométrica $r$--$T$ com curvas de UR & Aula 1, após variáveis \\
3 & LCL exato × regra; base de nuvem ao longo do dia & Aula 2, após LCL \\
4 & Água precipitável; escala de 7\,\%/K & fim da Aula 2 \\
\hline
\end{tabular}
\end{center}

O painel gelo × líquido do Caso~1 planta a semente do processo de Bergeron
(Cap.~7); o Caso~4 fecha com a conexão climática (chuvas extremas).

\subsection*{Exercícios complementares e soluções}

Nas notas: T1--T5 (integração de CC com $L_v$ variável, regra dos 7\,\%/K,
dedução de $r(\varepsilon)$, dedução da regra do LCL, água precipitável
analítica) e N1--N4 (células no notebook). Soluções em
\texttt{cap04\_solucoes.ipynb} (\emph{não distribuir}): T1 e T4 com
\texttt{sympy}, N2 com \texttt{brentq} validado contra o MetPy.
Pares teoria--numérica: T1$\leftrightarrow$N1 e T2$\leftrightarrow$N4.

\subsection*{Conexões com o resto do curso}

$e_s$ sobre gelo × líquido $\to$ Bergeron (Cap.~7); LCL $\to$ base de nuvem
e diagramas (Caps.~5--6); $r$ conservada + $\theta$ $\to$ diagramas
termodinâmicos (Cap.~5); água precipitável $\to$ rios atmosféricos e chuvas
extremas (Caps.~13, 21); balanço de água $\to$ ciclo hidrológico (Cap.~21).

\clearpage
% Guia do professor — Capítulo 5: Estabilidade Atmosférica
\section{Capítulo 5 --- Estabilidade Atmosférica}

\subsection*{Visão geral e ritmo}

O capítulo-síntese da primeira metade do curso: Caps.~1--4 se encontram nos
diagramas termodinâmicos e no método da parcela. Empuxo, Brunt--Väisälä,
estabilidade estática (com o método \emph{não local} que é marca registrada
do Stull) e número de Richardson. \textbf{Ritmo sugerido: 2 aulas de 2\,h}
(o capítulo é longo no livro; estas notas destilam o essencial).

\begin{itemize}
  \item \textbf{Aula 1} — diagramas termodinâmicos (construção passo a
        passo no quadro), método da parcela, empuxo, dedução do oscilador
        de Brunt--Väisälä. Fechar com Casos~1--2 (15\,min).
  \item \textbf{Aula 2} — estabilidade local × não local (o coração
        pedagógico da aula), Richardson e turbulência, leitura de perfis
        (tropopausa, camada de mistura, inversões). Casos~3--4 (15\,min).
\end{itemize}

\subsection*{Objetivos de aprendizagem}

\begin{enumerate}
  \item ler e usar um Skew-$T$: adiabáticas secas/úmidas, isoumas, LCL,
        trajetória de parcela;
  \item calcular o empuxo via temperatura virtual e estimar velocidades de
        termais;
  \item deduzir $N^2$ e interpretar o período de Brunt--Väisälä em
        fenômenos visíveis;
  \item aplicar o método não local e explicar por que o critério local
        falha;
  \item calcular $Ri$ e prever turbulência de céu claro;
  \item identificar tropopausa, camada de mistura e inversões num perfil.
\end{enumerate}

\subsection*{Pré-requisitos a reativar}

$T_v$ (Cap.~1), $\theta$ e $\Gamma_d$ (Cap.~3), LCL e $r_s$ (Cap.~4),
oscilador harmônico (física básica — os alunos vão sorrir ao reencontrá-lo
na atmosfera).

\subsection*{Armadilhas conceituais frequentes}

\begin{itemize}
  \item \textbf{``Gradiente ambiental $6{,}5 <$ adiabático $9{,}8$ $\Rightarrow$
        sempre estável''}: verdade para o ar \emph{não saturado} em média —
        mas o método da parcela compara com a \emph{adiabática úmida} acima
        do LCL, e a instabilidade condicional aparece. Preparar bem: é a
        ponte para o Cap.~14.
  \item \textbf{Estabilidade só pelo gradiente local}: a armadilha central;
        o exemplo da tarde convectiva desfaz.
  \item \textbf{Confundir $N$ com frequência de oscilação visível}: nuvens
        em faixa têm espaçamento espacial $U\cdot P_{BV}$; a conversão
        tempo$\to$espaço confunde.
  \item \textbf{$Ri$ como critério simétrico}: a transição
        laminar$\to$turbulento ($Ri<0{,}25$) e a volta ($Ri>1$) não
        coincidem — histerese.
  \item \textbf{Esquecer $T_v$ no empuxo}: em ar úmido tropical o erro
        chega a alguns décimos de K — exatamente a ordem do empuxo de um
        termal.
\end{itemize}

\subsection*{Demonstração de baixo custo}

Garrafa PET com água + azeite (ou água quente colorida embaixo de fria):
oscilações de interface ilustram Brunt--Väisälä; agitar mostra a transição
para turbulência (Richardson qualitativo). Vídeo de ondas de
Kelvin--Helmholtz em nuvens fecha a Aula 2 com aplauso.
\textbf{Material de apoio}: o \S5.11 do livro traz diagramas
termodinâmicos em página inteira (emagrama, Skew-$T$, tefigrama etc.)
prontos para imprimir — distribuir cópias para o trabalho manual em
sala e nas provas com consulta.

\subsection*{Problemas sugeridos do livro (exercícios do Cap.~5)}

Aquecimento: uso do diagrama termodinâmico (plotar sondagem, achar LCL);
cálculo de $N$ e período; classificação de estabilidade de perfis dados.
Aprofundamento: método não local em perfis com inversão; $Ri$ com
cisalhamento realista. Síntese: ``por que o topo da camada de mistura
coincide com a base das nuvens em dias de bom tempo?''.
\emph{Confira a numeração na sua cópia (v1.02b).}

\subsection*{Uso do notebook \texttt{cap05\_estabilidade.ipynb}}

\begin{center}
\begin{tabular}{lll}
\hline
Caso & Conteúdo & Momento sugerido \\
\hline
1 & Diagrama do zero + Skew-$T$ com parcela e área de empuxo & Aula 1, após diagramas \\
2 & Oscilador de Brunt--Väisälä (EDO + amortecimento) & Aula 1, após $N^2$ \\
3 & Estabilidade não local implementada & Aula 2, após o método \\
4 & Perfil de $Ri$ e camadas de CAT & Aula 2, após Richardson \\
\hline
\end{tabular}
\end{center}

O Caso~1 tem intenção dupla: desmistificar o diagrama (construindo-o) e
plantar o CAPE (área sombreada) para o Cap.~14. O Caso~3 é a implementação
fiel do algoritmo do livro — raro em qualquer curso.

\subsection*{Exercícios complementares e soluções}

Nas notas: T1--T5 (empuxo com \emph{liquid loading}, dedução de $N^2$,
períodos por camada, CAT no jato, contraexemplo local × não local) e
N1--N4. Soluções em \texttt{cap05\_solucoes.ipynb} (\emph{não distribuir}):
T2 com \texttt{sympy} (solução do oscilador), N1 com amortecimento, N3 com
o algoritmo comentado linha a linha. Pares: T2$\leftrightarrow$N1,
T5$\leftrightarrow$N3.

\subsection*{Conexões com o resto do curso}

Diagramas + parcela $\to$ CAPE e tempestades (Cap.~14); não local $\to$
camada limite (Cap.~18); inversões $\to$ dispersão de poluentes (Cap.~19);
$N$ $\to$ ondas de montanha (Cap.~17); $Ri$ $\to$ turbulência e aviação.

\clearpage
% Guia do professor — Capítulo 6: Nuvens
\section{Capítulo 6 --- Nuvens}

\subsection*{Visão geral e ritmo}

O capítulo mais visual do curso: a taxonomia das nuvens (linguagem
operacional) e os três mecanismos de saturação (física), com nevoeiros como
estudo de caso e a geometria fractal como sobremesa para físicos.
\textbf{Ritmo sugerido: 2 aulas de 2\,h} — pode comprimir para 1 aula e
meia se o semestre apertar (a classificação anda rápido com fotos).

\begin{itemize}
  \item \textbf{Aula 1} — gêneros com fotos (15\,min bem gastos),
        gramática dos nomes, os três caminhos para saturar (com a dedução
        da convexidade), bafo e contrails. Caso~1 (10\,min).
  \item \textbf{Aula 2} — nevoeiros (tabela + EDO do nevoeiro de
        radiação), organização (ruas, ondas, células), fractais, oktas.
        Casos~2--4 (20\,min; o Caso~3 rende).
\end{itemize}

\subsection*{Objetivos de aprendizagem}

\begin{enumerate}
  \item nomear os 10 gêneros e inferir o mecanismo (estável × convectivo)
        do aspecto;
  \item explicar os três caminhos para a saturação na curva $e_s(T)$ e
        provar por que mistura satura (convexidade);
  \item classificar nevoeiros pelo mecanismo e estimar tempo de formação do
        nevoeiro de radiação;
  \item relacionar padrões de organização com $z_i$ e com $N$ (Cap.~5);
  \item explicar o que significa dimensão fractal $D \simeq 1{,}35$ e como
        se mede;
  \item usar oktas e teto na linguagem METAR.
\end{enumerate}

\subsection*{Pré-requisitos a reativar}

Curva $e_s(T)$ e LCL (Cap.~4); $N$ de Brunt--Väisälä e camada de mistura
(Cap.~5); balanço radiativo noturno (Cap.~2). Fractais: nenhum
pré-requisito — introduzir do zero é parte da graça.

\subsection*{Armadilhas conceituais frequentes}

\begin{itemize}
  \item \textbf{``Nuvem é vapor''}: nuvem é água \emph{líquida/gelo};
        vapor é invisível. Corrigir logo com o bafo.
  \item \textbf{``Misturar ar só faz média''}: a média é na corda, mas a
        saturação é na curva — a convexidade surpreende até veteranos.
  \item \textbf{Nevoeiro como ``nuvem baixa'' sem física}: cada tipo tem um
        mecanismo distinto; a tabela de 5 tipos organiza.
  \item \textbf{Contrail como ``fumaça''}: é nuvem de mistura; a fuligem só
        fornece núcleos (gancho para o Cap.~7).
  \item \textbf{Perímetro fractal}: alunos esperam convergência ao refinar
        a régua; a divergência $\ell^{1-D}$ é contraintuitiva — o exercício
        T4 fixa.
\end{itemize}

\subsection*{Demonstração de baixo custo}

Nuvem na garrafa: garrafa PET com um dedo d'água + fumaça de palito de
fósforo; apertar e soltar — expansão adiabática forma a nuvem no ato
(mecanismo 2 + núcleos do Cap.~7). Complemento: vídeo timelapse de nevoeiro
de radiação se formando em vale.

\subsection*{Problemas sugeridos do livro (exercícios do Cap.~6)}

Aquecimento: identificação de gêneros a partir de descrições; altura de
base de cumulus via LCL. Aprofundamento: problemas de mistura
(bafo/contrail); tempo de formação de nevoeiro. Síntese: ``por que noites
nubladas não dão nevoeiro de radiação?''.
\emph{Confira a numeração na sua cópia (v1.02b).}

\subsection*{Uso do notebook \texttt{cap06\_nuvens.ipynb}}

\begin{center}
\begin{tabular}{lll}
\hline
Caso & Conteúdo & Momento sugerido \\
\hline
1 & Diagrama de mistura: bafo e contrail em $e_s(T)$ & Aula 1, após mistura \\
2 & EDO do nevoeiro de radiação (forma e dissipa) & Aula 2, após nevoeiros \\
3 & Nuvens fractais sintetizadas; $D$ por contagem de caixas & Aula 2, fractais \\
4 & $\lambda$ de ondas e ruas a partir da sondagem & Aula 2, organização \\
\hline
\end{tabular}
\end{center}

O Caso~3 é o mais ``físico computacional'' do curso até aqui (síntese
espectral + contagem de caixas) — bom para apresentar como mini-projeto.

\subsection*{Exercícios complementares e soluções}

Nas notas: T1--T5 (convexidade e supersaturação de misturas, critério de
Appleman, dedução do $t_{fog}$, aritmética fractal, espaçamentos de
organização) e N1--N4. Soluções em \texttt{cap06\_solucoes.ipynb}
(\emph{não distribuir}): T1 com \texttt{sympy} (segunda derivada de
Clausius--Clapeyron), N3 com o estudo $D(\beta)$ completo. Pares:
T1$\leftrightarrow$N1, T3$\leftrightarrow$N2, T4$\leftrightarrow$N3.

\subsection*{Conexões com o resto do curso}

LCL/base de cumulus $\to$ Cap.~4; estável × convectivo e $N$ $\to$ Cap.~5;
núcleos de condensação e crescimento de gotas $\to$ Cap.~7; oktas/teto
$\to$ METARs (Cap.~9); nevoeiro de radiação $\to$ camada limite noturna
(Cap.~18); células abertas/fechadas $\to$ interação ar-mar (Cap.~21).

\clearpage
% Guia do professor — Capítulo 7: Processos de Precipitação
\section{Capítulo 7 --- Processos de Precipitação}

\subsection*{Visão geral e ritmo}

Microfísica de nuvens organizada como um problema de escalas: por que a
natureza precisa de aerossol para condensar e de colisões (ou gelo) para
chover. \textbf{Ritmo sugerido: 2 aulas de 2\,h.}

\begin{itemize}
  \item \textbf{Aula 1} — escada de tamanhos, Kelvin, Köhler, ativação,
        nucleação de gelo. Casos~1--2 (15\,min; o \texttt{pyrcel} rende).
  \item \textbf{Aula 2} — difusão ($\sqrt t$ e seu gargalo),
        colisão--coalescência, Bergeron, velocidades terminais,
        Marshall--Palmer. Casos~3--4 (15\,min).
\end{itemize}

\subsection*{Objetivos de aprendizagem}

\begin{enumerate}
  \item explicar por que a atmosfera não condensa sem CCN (barreira de
        Kelvin);
  \item ler uma curva de Köhler: $r^*$, $S^*$, ativação e papel do núcleo;
  \item mostrar que difusão dá $r\propto\sqrt t$ e por que isso não faz
        chuva;
  \item comparar chuva quente (coleta) e fria (Bergeron) e saber qual domina
        onde;
  \item usar Marshall--Palmer para calcular massa, taxa de chuva e
        refletividade.
\end{enumerate}

\subsection*{Pré-requisitos a reativar}

$e_s$ sobre água/gelo (Cap.~4); tensão superficial e energia livre
(termodinâmica básica); arrasto de Stokes (mecânica dos fluidos elementar).

\subsection*{Armadilhas conceituais frequentes}

\begin{itemize}
  \item \textbf{``Nuvem cresce até virar chuva por condensação''}: o
        gargalo do $\sqrt t$ desfaz; insistir nos tempos (minutos vs.\
        dias).
  \item \textbf{Supersaturações imaginárias}: alunos assumem UR de
        110--120\,\% em nuvens; o real é $\lesssim 100{,}5$\,\% — e é por
        isso que Kelvin é barreira.
  \item \textbf{Gelo a $0\,°$C}: água super-resfriada até $-38\,°$C
        surpreende; núcleos de gelo são raros.
  \item \textbf{Marshall--Palmer como lei universal}: é climatologia de
        chuva estratiforme; convecção e neve fogem dela.
  \item \textbf{Efeito do aerossol}: mais CCN $=$ mais gotas \emph{menores}
        (não mais chuva!) — nuvem mais brilhante e chuva possivelmente
        suprimida.
\end{itemize}

\subsection*{Demonstração de baixo custo}

Garrafa PET (nuvem na garrafa) com e sem fumaça: sem núcleos quase nada
acontece; com fumaça a nuvem é instantânea — Kelvin e CCN em 2 minutos.
Vídeo de câmara de nuvens (cloud chamber) fecha a Aula 1.

\subsection*{Problemas sugeridos do livro (exercícios do Cap.~7)}

Aquecimento: UR de equilíbrio de gotas pequenas; $r^*$ e $S^*$ para núcleos
dados; velocidades terminais. Aprofundamento: tempos de crescimento por
difusão; taxa de chuva de distribuições dadas. Síntese: ``por que
semeadura de nuvens usa AgI?'' (estrutura cristalina $\approx$ gelo).
\emph{Confira a numeração na sua cópia (v1.02b).}

\subsection*{Uso do notebook \texttt{cap07\_precipitacao.ipynb}}

\begin{center}
\begin{tabular}{lll}
\hline
Caso & Conteúdo & Momento sugerido \\
\hline
1 & Curvas de Köhler; $r^*(r_d)$, $S^*(r_d)$ & Aula 1, após Köhler \\
2 & \texttt{pyrcel}: ativação vs.\ updraft & fim da Aula 1 \\
3 & Difusão × coleta: o gargalo e a avalanche & Aula 2, após crescimento \\
4 & Marshall--Palmer: momentos, $R$, $Z$, relação $Z$--$R$ & fim da Aula 2 \\
\hline
\end{tabular}
\end{center}

O Caso~2 usa um pacote de pesquisa de verdade (\texttt{pyrcel}) — vale
contar aos alunos que é o mesmo tipo de módulo que roda dentro de modelos
climáticos. O Caso~4 fecha preparando o radar do Cap.~8.

\subsection*{Exercícios complementares e soluções}

Nas notas: T1--T5 (Kelvin por energia livre, extremos de Köhler, tempos de
difusão, Stokes, momentos de M--P) e N1--N4. Soluções em
\texttt{cap07\_solucoes.ipynb} (\emph{não distribuir}): T2 e T5 com
\texttt{sympy} (extremos e integrais $\Gamma$), N4 verificando $Z = aR^b$.
Pares: T2$\leftrightarrow$N1, T5$\leftrightarrow$N4.

\subsection*{Conexões com o resto do curso}

$e_s$ água/gelo $\to$ Cap.~4; correntes ascendentes que sustentam gotas
$\to$ Caps.~5 e 14; $Z$--$R$ $\to$ radar (Cap.~8); aerossol--nuvem $\to$
clima (Cap.~21); granizo $\to$ Cap.~15.

\clearpage
% Guia do professor — Capítulo 8: Satélites e Radar
\section{Capítulo 8 --- Satélites e Radar}

\subsection*{Visão geral e ritmo}

O Cap.~2 vestido de operacional: órbitas por Newton, funções peso como
coração da sondagem por satélite, e o radar do $D^6$ ao dilema do Doppler.
\textbf{Ritmo sugerido: 2 aulas de 2\,h} (1 aula e meia se necessário,
cortando detalhes de imagem).

\begin{itemize}
  \item \textbf{Aula 1} — órbitas (dedução geoestacionária no quadro),
        funções peso (a dedução central), canais Vis/IV/vapor com imagens
        do dia. Casos~1--2 (15\,min).
  \item \textbf{Aula 2} — radar: equação, $Z$, $Z$--$R$ e armadilhas;
        Doppler e aliasing. Casos~3--4 (15\,min).
\end{itemize}

\subsection*{Objetivos de aprendizagem}

\begin{enumerate}
  \item deduzir a altitude geoestacionária e comparar geo × polar;
  \item explicar a função peso e por que o pico está em $\tau = 1$;
  \item interpretar os canais visível, IV e vapor d'água (e combinações);
  \item converter dBZ em chuva com consciência das incertezas;
  \item explicar o dilema do Doppler e reconhecer aliasing.
\end{enumerate}

\subsection*{Pré-requisitos a reativar}

Gravitação circular (física 1); Planck/Stefan--Boltzmann e Beer (Cap.~2);
Marshall--Palmer e o $D^6$ (Cap.~7); atmosfera padrão para converter
temperatura de topo em altura (Cap.~1).

\subsection*{Armadilhas conceituais frequentes}

\begin{itemize}
  \item \textbf{``Satélite vê a superfície''}: no IV ele vê o nível de
        $\tau\simeq1$ do canal — que pode ser vapor a 8\,km.
  \item \textbf{IV de cirrus fino}: semitransparente mistura chão quente e
        nuvem fria — a altura inferida sai baixa demais (T3).
  \item \textbf{dBZ como chuva exata}: $Z$--$R$ é estatística; a
        propagação de $\pm3$\,dB (T5) cura a fé.
  \item \textbf{Doppler mede só a componente radial}: vento perpendicular
        ao feixe é invisível.
  \item \textbf{Aliasing}: velocidades ``impossíveis'' coladas a um jato
        forte costumam ser dobras, não fenômenos.
\end{itemize}

\subsection*{Demonstração de baixo custo}

Imagens em tempo real (GOES-East no site da NOAA/CPTEC) nos três canais da
mesma cena — classificar nuvens ao vivo com a tabela da aula. Doppler:
vídeo de radar do Redemet/simulação de aliasing do notebook.

\subsection*{Problemas sugeridos do livro (exercícios do Cap.~8)}

Aquecimento: períodos orbitais; conversões dBZ$\leftrightarrow$$R$; alturas
de topo pelo IV. Aprofundamento: funções peso para perfis dados; $V_{max}$
e $R_{max}$ para PRFs reais. Síntese: ``por que não existe um único radar
perfeito?'' (dilema + atenuação + custo).
\emph{Confira a numeração na sua cópia (v1.02b).}

\subsection*{Uso do notebook \texttt{cap08\_sensoriamento.ipynb}}

\begin{center}
\begin{tabular}{lll}
\hline
Caso & Conteúdo & Momento sugerido \\
\hline
1 & Órbitas: $T(z)$, cobertura & Aula 1, após Kepler \\
2 & Funções peso: sondador de 6 canais & Aula 1, após dedução \\
3 & $Z$--$R$: inversão e incerteza & Aula 2, após radar \\
4 & Aliasing Doppler e dealias & fim da Aula 2 \\
\hline
\end{tabular}
\end{center}

Projeto opcional para equipes: ler um volume real de radar com
\texttt{Py-ART} (instalado no ambiente) ou uma cena de satélite com
\texttt{Satpy} — ambos têm tutoriais oficiais excelentes.

\subsection*{Exercícios complementares e soluções}

Nas notas: T1--T5 e N1--N4. Soluções em \texttt{cap08\_solucoes.ipynb}
(\emph{não distribuir}); T2 com \texttt{sympy} (máximo da função peso),
N4 com o dealias por continuidade e seus modos de falha.
Pares: T2$\leftrightarrow$N2, T4$\leftrightarrow$N4.

\subsection*{Conexões com o resto do curso}

Beer/Planck $\to$ Cap.~2; $D^6$/M--P $\to$ Cap.~7; imagens e cartas $\to$
Cap.~9; mesociclone no Doppler $\to$ Cap.~15; assimilação de radiâncias
$\to$ Cap.~20.

\clearpage
% Guia do professor — Capítulo 9: Boletins e Análise de Cartas
\section{Capítulo 9 --- Boletins Meteorológicos e Análise de Cartas}

\subsection*{Visão geral e ritmo}

Capítulo de alfabetização operacional: METAR, modelo de estação, redução ao
nível do mar e análise (manual e objetiva). \textbf{Ritmo: 1 aula
expositiva + 1 aula prática} com carta impressa e lápis — a prática é o
coração.

\begin{itemize}
  \item \textbf{Aula 1} — METAR grupo a grupo (com boletins do dia!),
        modelo de estação, redução ao nível do mar. Casos~1--2 (10\,min).
  \item \textbf{Aula 2 (prática)} — cada dupla analisa a mesma carta de
        superfície impressa: isóbaras de 4 em 4\,hPa, centros, frentes
        tentativas. Fechar comparando com a análise oficial e com o
        Caso~3.
\end{itemize}

\subsection*{Objetivos de aprendizagem}

\begin{enumerate}
  \item decodificar METARs sem tabela de consulta (grupos principais);
  \item plotar e ler o modelo de estação (barbela, pressão codificada);
  \item explicar e aplicar a redução ao nível do mar;
  \item traçar isóbaras seguindo as regras e identificar centros;
  \item descrever a análise objetiva (Cressman/Barnes) e seu papel
        histórico.
\end{enumerate}

\subsection*{Pré-requisitos a reativar}

Hipsométrica (Cap.~1); oktas (Cap.~6); $T_d$/UR (Cap.~4); convenção de
direção do vento (Cap.~1 — o erro clássico reaparece na barbela).

\subsection*{Armadilhas conceituais frequentes}

\begin{itemize}
  \item \textbf{Barbela ao contrário}: a haste aponta de onde o vento vem;
        metade da turma erra na primeira vez.
  \item \textbf{Pressão codificada}: 052 $\to$ 1005{,}2 (não 905{,}2);
        a regra do valor mais próximo de 1000 e sua exceção em furacões.
  \item \textbf{QNH × QFE × pressão da estação}: três pressões diferentes;
        cartas usam a reduzida.
  \item \textbf{Isóbaras ``ligando pontos''}: interpolação proporcional,
        não conexão de estações.
  \item \textbf{Análise objetiva como verdade}: é filtro; o raio de
        influência decide o que sobrevive (T5/N3).
\end{itemize}

\subsection*{Demonstração de baixo custo}

METARs ao vivo do REDEMET/aviationweather.gov no projetor: decodificar o
aeroporto da cidade em tempo real. Para a prática: imprimir a carta de
superfície da Marinha/CPTEC do dia anterior sem análise.

\subsection*{Problemas sugeridos do livro (exercícios do Cap.~9)}

Aquecimento: decodificação de boletins; plotagem de estações.
Aprofundamento: análise de mini-cartas com 10--15 estações. Síntese:
``que observações faltam sobre os oceanos e como os satélites (Cap.~8)
preenchem?''. \emph{Confira a numeração na sua cópia (v1.02b).}

\subsection*{Uso do notebook \texttt{cap09\_cartas.ipynb}}

\begin{center}
\begin{tabular}{lll}
\hline
Caso & Conteúdo & Momento sugerido \\
\hline
1 & Parser de METAR (Brasil) & Aula 1, após o código \\
2 & Modelo de estação (MetPy) & Aula 1, após o modelo \\
3 & Cressman: pontos $\to$ isóbaras & Aula 2, após a prática manual \\
4 & Redução ao nível do mar: antes × depois & Aula 1 ou 2 \\
\hline
\end{tabular}
\end{center}

Momento pedagógico forte: mostrar o Caso~3 \emph{depois} da análise manual
— os alunos reconhecem no algoritmo o que a mão deles acabou de fazer.

\subsection*{Exercícios complementares e soluções}

Nas notas: T1--T5 e N1--N4 (o par N1$\to$N2 constrói um mini-sistema de
plotagem). Soluções em \texttt{cap09\_solucoes.ipynb}
(\emph{não distribuir}); N3 traz validação cruzada de verdade — primeira
aparição de um conceito de aprendizado estatístico no curso.

\subsection*{Conexões com o resto do curso}

Hipsométrica $\to$ Cap.~1; oktas/teto $\to$ Cap.~6; isóbaras$\to$vento
$\to$ Cap.~10; frentes $\to$ Cap.~12; análise objetiva $\to$ assimilação
(Cap.~20).

\clearpage
% Guia do professor — Capítulo 10: Forças e Ventos
\section{Capítulo 10 --- Forças Atmosféricas e Ventos}

\subsection*{Visão geral e ritmo}

A porta de entrada da dinâmica: Coriolis, equilíbrio geostrófico, vento
gradiente, atrito/Ekman e continuidade — terminando na frase que organiza a
sinótica: \emph{baixas convergem e chovem; altas divergem e abrem}.
\textbf{Ritmo: 2 aulas de 2\,h.} Para físicos, é o capítulo mais
confortável do curso — aproveite para ir fundo nas deduções.

\begin{itemize}
  \item \textbf{Aula 1} — forças, Coriolis (com a desmistificação da pia),
        círculos inerciais, dedução geostrófica completa, Buys-Ballot,
        régua na carta. Casos~1--2 (15\,min).
  \item \textbf{Aula 2} — vento gradiente (quadrática no quadro), 3 forças
        com atrito e o desenho convergência/divergência, Ekman,
        continuidade e $w$. Casos~3--4 (15\,min).
\end{itemize}

\subsection*{Objetivos de aprendizagem}

\begin{enumerate}
  \item escrever a equação do movimento horizontal e a escala de cada
        termo;
  \item deduzir o vento geostrófico e usá-lo como régua de carta
        (Buys-Ballot no HS);
  \item explicar sub/supergeostrofia em baixas/altas e o regime
        ciclostrófico;
  \item desenhar o equilíbrio de 3 forças e concluir
        convergência$\to$subida em baixas;
  \item derivar/verificar a espiral de Ekman e estimar $w$ por
        continuidade.
\end{enumerate}

\subsection*{Pré-requisitos a reativar}

Referenciais não inerciais e força centrífuga/Coriolis (mecânica);
gradiente e divergência (cálculo vetorial); hidrostática (Cap.~1) para a
versão em coordenadas de pressão (mencionar, não exigir).

\subsection*{Armadilhas conceituais frequentes}

\begin{itemize}
  \item \textbf{A pia}: Coriolis em escala de pia é $10^{-7}$ da força
        real; medir com o T1 (escala de tempo).
  \item \textbf{HS × HN}: todo desenho de livro americano gira ao
        contrário aqui; treinar Buys-Ballot no HS até virar reflexo.
  \item \textbf{``Coriolis faz o vento''}: Coriolis só \emph{desvia}; quem
        empurra é a PGF.
  \item \textbf{Geostrofia como lei}: é equilíbrio assintótico; perto do
        equador, em curvas fechadas e na CL ele falha — e as falhas são o
        conteúdo dos Caps.~11--16.
  \item \textbf{Sinal de $f$}: erros de sinal no HS são epidêmicos em
        listas; cobrar o desenho junto com a conta.
\end{itemize}

\subsection*{Demonstração de baixo custo}

Globo girante (ou cadeira giratória) + bolinha rolada: o desvio aparente
ao filmar do referencial girante. Vídeo clássico do carrossel do MIT
fecha a discussão da pia.

\subsection*{Problemas sugeridos do livro (exercícios do Cap.~10)}

Aquecimento: $f$ em várias latitudes; $U_g$ de cartas dadas; período
inercial. Aprofundamento: vento gradiente em baixas/altas; ângulo de
cruzamento com atrito. Síntese: ``por que o centro de uma alta é sempre de
ventos fracos?'' (T3 responde).
\emph{Confira a numeração na sua cópia (v1.02b).}

\subsection*{Uso do notebook \texttt{cap10\_ventos.ipynb}}

\begin{center}
\begin{tabular}{lll}
\hline
Caso & Conteúdo & Momento sugerido \\
\hline
1 & Círculos inerciais e cicloides & Aula 1, após Coriolis \\
2 & Geostrófico × gradiente na baixa sintética & Aula 1/2, após dedução \\
3 & Ekman: analítica × numérica & Aula 2, após atrito \\
4 & Convergência $\to$ mapa de $w$ & fim da Aula 2 \\
\hline
\end{tabular}
\end{center}

O Caso~4 é o fecho conceitual: o mapa de $w$ da baixa sintética É o mapa
de nuvens que os alunos verão nas imagens de satélite do Cap.~13.

\subsection*{Exercícios complementares e soluções}

Nas notas: T1--T5 (inerciais, equador, teto das altas, ângulo de
cruzamento, verificação de Ekman + transporte) e N1--N4. Soluções em
\texttt{cap10\_solucoes.ipynb} (\emph{não distribuir}): T3 e T5 com
\texttt{sympy}. Pares: T1$\leftrightarrow$N1, T5$\leftrightarrow$N3.

\subsection*{Conexões com o resto do curso}

Isóbaras$\to$vento fecha o Cap.~9; convergência$\to w$ abre os
Caps.~12--13; Ekman e cruzamento $\to$ camada limite (Cap.~18) e furacões
(Cap.~16); ciclostrófico $\to$ tornados (Cap.~15); geostrofia em altitude
$\to$ vento térmico (Cap.~11).

\clearpage
% Guia do professor — Capítulo 11: Circulação Geral
\section{Capítulo 11 --- Circulação Geral}

\subsection*{Visão geral e ritmo}

O capítulo que conecta o micro (Caps.~1--10) ao planeta: desequilíbrio
radiativo, células, jatos (momento angular e vento térmico) e ondas de
Rossby. O capítulo do livro é longo (60\,pp); estas notas destilam a
espinha dorsal física. \textbf{Ritmo: 2 aulas de 2\,h.}

\begin{itemize}
  \item \textbf{Aula 1} — os dois perfis radiativos e o transporte exigido
        (integral no quadro), Hadley, momento angular e o jato subtropical,
        o mapa-múndi climático. Casos~1--2 (15\,min; o EBM rende).
  \item \textbf{Aula 2} — vento térmico (dedução completa), jato polar,
        Rossby ($\beta$ como restauração), cinturões e bloqueios.
        Casos~3--4 (15\,min).
\end{itemize}

\subsection*{Objetivos de aprendizagem}

\begin{enumerate}
  \item calcular o transporte meridional exigido a partir dos perfis
        radiativos;
  \item explicar Hadley, seus limites ($\sim$30°) e o jato subtropical por
        momento angular;
  \item deduzir e aplicar o vento térmico; localizar o jato polar;
  \item explicar a restauração $\beta$ e usar $c = U - \beta/k^2$;
  \item ler o mapa-múndi climático como consequência das células.
\end{enumerate}

\subsection*{Pré-requisitos a reativar}

Insolação por latitude (Cap.~2, Caso~2!); geostrofia e hidrostática
(Caps.~1, 10); momento angular (mecânica); vorticidade (definir
rapidamente — físicos aceitam $\zeta = \nabla\times\vec U$ de imediato).

\subsection*{Armadilhas conceituais frequentes}

\begin{itemize}
  \item \textbf{``Hadley vai até o polo''}: o modelo de uma célula falha
        por momento angular (134\,m/s!) e baroclinia; T2 quantifica.
  \item \textbf{Jato como ``rio de ar'' isolado}: é a expressão vertical do
        contraste térmico (vento térmico), não uma entidade à parte.
  \item \textbf{Rossby para leste}: a FASE das ondas longas anda para
        oeste em relação ao escoamento; a advecção pelo $U$ costuma vencer
        — distinguir fase × grupo × advecção.
  \item \textbf{ZCIT no equador geográfico}: migra com a estação e fica em
        média ao norte — vale mostrar a climatologia.
  \item \textbf{Sinais no HS}: vento térmico e Buys-Ballot de novo com
        $f<0$; insistir nos desenhos com o sul para baixo.
\end{itemize}

\subsection*{Demonstração de baixo custo}

Tanque girante improvisado (balde em prato giratório com água + corante +
gelo no centro): células e ondas aparecem. Alternativa: vídeos do
laboratório de tanque do MIT (Marshall) — Hadley e baroclinia em 3\,min.

\subsection*{Problemas sugeridos do livro (exercícios do Cap.~11)}

Aquecimento: $\beta$ e período de Rossby em várias latitudes; vento térmico
com gradientes dados. Aprofundamento: transporte exigido de tabelas
radiativas; jato por camadas. Síntese: ``o que acontece com os cinturões
num planeta girando 2$\times$ mais rápido?''.
\emph{Confira a numeração na sua cópia (v1.02b).}

\subsection*{Uso do notebook \texttt{cap11\_circulacao.ipynb}}

\begin{center}
\begin{tabular}{lll}
\hline
Caso & Conteúdo & Momento sugerido \\
\hline
1 & Transporte exigido (os 5\,PW) & Aula 1, após a integral \\
2 & EBM meridional (\texttt{climlab}): $T(\phi)$, linha do gelo & Aula 1 \\
3 & Dispersão de Rossby e onda estacionária & Aula 2, após Rossby \\
4 & Vento térmico $\to$ jato num perfil realista & Aula 2, após dedução \\
\hline
\end{tabular}
\end{center}

O Caso~2 reaproveita o \texttt{climlab} do Cap.~2 em modo meridional — e
antecipa o gelo-albedo do Cap.~21.

\subsection*{Exercícios complementares e soluções}

Nas notas: T1--T5 e N1--N4. Soluções em \texttt{cap11\_solucoes.ipynb}
(\emph{não distribuir}): T3 com \texttt{sympy}; N2 é a ponte direta para o
Cap.~21. Pares: T1$\leftrightarrow$N1, T5$\leftrightarrow$N3.

\subsection*{Conexões com o resto do curso}

Insolação (Cap.~2); geostrofia/vorticidade (Cap.~10); ciclones nascem nas
ondas (Cap.~13); jato e CAT (Cap.~5); bloqueios e extremos (Cap.~21);
correntes oceânicas e ENSO (Cap.~21).

\clearpage
% Guia do professor — Capítulo 12: Massas de Ar e Frentes
\section{Capítulo 12 --- Massas de Ar e Frentes}

\subsection*{Visão geral e ritmo}

A ponte entre a circulação geral (Cap.~11) e os ciclones (Cap.~13): onde
o ar ganha identidade (massas) e onde as identidades colidem (frentes).
Capítulo muito visual — cortes verticais e meteogramas rendem mais que
equações, com duas exceções nobres: Margules e a frontogênese
exponencial. \textbf{Ritmo: 2 aulas de 2\,h.}

\begin{itemize}
  \item \textbf{Aula 1} — anticiclones como fábricas, classificação de
        Bergeron com $(\theta_w, r)$, gênese fria (EDO no quadro),
        modificação em trânsito. Caso~1 (10\,min).
  \item \textbf{Aula 2} — anatomia comparada fria × quente, Margules,
        frontogênese por deformação, oclusões e linha seca. Casos~2--4
        (20\,min; o meteograma do Caso~4 fecha bem a aula).
\end{itemize}

\subsection*{Objetivos de aprendizagem}

\begin{enumerate}
  \item classificar massas de ar por origem e pelo par conservado
        $(\theta_w, r)$;
  \item estimar o tempo de gênese fria pela EDO radiativa linearizada;
  \item calcular a inclinação frontal por Margules e ligá-la à
        nebulosidade de cada frente;
  \item explicar a frontogênese por deformação e sua taxa exponencial;
  \item reconhecer passagens frontais em meteogramas e cartas
        (isóbaras em ``V'', giro de vento no HS).
\end{enumerate}

\subsection*{Pré-requisitos a reativar}

$\theta_w$ e $r$ como etiquetas conservadas (Cap.~4); fluxos bulk
(Cap.~3); geostrofia e hidrostática (Caps.~1, 10); Stefan--Boltzmann
(Cap.~2); leitura de cartas (Cap.~9).

\subsection*{Armadilhas conceituais frequentes}

\begin{itemize}
  \item \textbf{``Frente é uma linha''}: é uma \emph{zona} 3-D com
        $\sim$100\,km de largura e uma cunha quase deitada (1/100!) —
        o desenho em escala real é chocante e vale fazer.
  \item \textbf{Giro de vento no HS}: N$\to$SW na fria (não S$\to$NW,
        que é o espelho do HN dos livros importados). Insistir com a
        rosa dos ventos no quadro.
  \item \textbf{``A frente fria é fria''}: o que define é o
        \emph{avanço} do ar frio; num pós-frontal de céu limpo a tarde
        pode ser mais quente que a manhã pré-frontal.
  \item \textbf{Margules sem cisalhamento}: sem $\Delta U$ ao longo da
        frente não há interface inclinada estacionária — frentes são
        estruturas dinâmicas, não ``paredes de densidade''.
  \item \textbf{Oclusão comparando $T$}: compara-se $\theta_w$ das
        massas frias (T5); com $T$ crua o critério erra sob inversões.
\end{itemize}

\subsection*{Demonstração de baixo custo}

Tanque de água com divisória: água fria tingida de um lado, quente do
outro; ao remover a divisória a cunha fria avança e a inclinação
aparece (sem rotação ela desaba — deixe explícito que é Margules SEM o
$f$, por isso não estaciona). Meteograma real: baixar a série da
estação do IAG/aeroporto na última friagem e sobrepor ao Caso~4.

\subsection*{Problemas sugeridos do livro (exercícios do Cap.~12)}

Aquecimento: classificação de massas para cidades dadas; tempo de
gênese com números diferentes. Aprofundamento: inclinação frontal com
dados de sondagem dupla; frontogênese com campo de deformação dado.
Síntese: ``por que o HS tem frentes mais zonais e menos violentas que o
HN?''. \emph{Confira a numeração na sua cópia (v1.02b).}

\subsection*{Uso do notebook \texttt{cap12\_frentes.ipynb}}

\begin{center}
\begin{tabular}{lll}
\hline
Caso & Conteúdo & Momento sugerido \\
\hline
1 & gênese cP: EDO radiativa ($T^4$) vs.\ linearizada & Aula 1 \\
2 & inclinação de Margules: varredura $\Delta T$, $\Delta U$ & Aula 2 \\
3 & frontogênese por deformação (mini-NWP, exata) & Aula 2 \\
4 & meteograma sintético da frente fria & fecho da Aula 2 \\
\hline
\end{tabular}
\end{center}

O Caso~3 é deliberadamente a ``previsão numérica em miniatura'' da vez:
advecção com solução exata — compare com a versão em grade (upwind) do
Cap.~3 e anuncie o Cap.~20.

\subsection*{Exercícios complementares e soluções}

Nas notas: T1--T5 e N1--N4. Soluções em \texttt{cap12\_solucoes.ipynb}
(\emph{não distribuir}): T1 e T2 verificados com \texttt{sympy}; N3
responde uma pergunta genuína (rotação não frontogeniza). Pares:
T2$\leftrightarrow$N1, T1$\leftrightarrow$N2, T3$\leftrightarrow$N3.

\subsection*{Conexões com o resto do curso}

Anticiclones e subsidência (Caps.~10--11); $(\theta_w, r)$ (Cap.~4);
nuvens frontais (Cap.~6); meteograma e cartas (Cap.~9); a frente que
ondula e fecha em ciclone é o Cap.~13 inteiro; linha seca e severo
(Cap.~14); ZCAS e chuvas persistentes (Cap.~21, variabilidade).

\clearpage
% Guia do professor — Capítulo 13: Ciclones Extratropicais
\section{Capítulo 13 --- Ciclones Extratropicais}

\subsection*{Visão geral e ritmo}

O capítulo mais dinâmico do curso: aqui a turma de física finalmente vê
mecânica dos fluidos geofísica de verdade (vorticidade, instabilidade,
realimentação). O fio narrativo: o ciclone é a máquina que faz o
transporte do Cap.~11, usando a energia da baroclinia via os mecanismos
do Cap.~10. \textbf{Ritmo: 2 aulas de 2\,h.}

\begin{itemize}
  \item \textbf{Aula 1} — ciclo de vida norueguês (4 desenhos), equação
        da vorticidade com foco no esticamento, advecção de vorticidade
        na onda (Caso~1), Eady e o porquê dos 4000\,km (Caso~2).
  \item \textbf{Aula 2} — EDO de spin-up (Caso~3), tendência de
        pressão, jet streaks (Caso~4), autodesenvolvimento e bombas.
\end{itemize}

\subsection*{Objetivos de aprendizagem}

\begin{enumerate}
  \item narrar o ciclo de vida norueguês com os tempos certos (dias);
  \item usar o termo de esticamento para explicar spin-up e spin-down;
  \item estimar taxa de crescimento e tamanho de ciclones com Eady;
  \item localizar regiões de ciclogênese por advecção de vorticidade e
        quadrantes de jet streak (com os sinais do HS);
  \item explicar a queda de pressão como exportação de massa da coluna.
\end{enumerate}

\subsection*{Pré-requisitos a reativar}

Vorticidade e continuidade (Cap.~10); vento térmico e Rossby (Cap.~11);
frentes e oclusão (Cap.~12); $N$ (Cap.~5); peso da coluna (Cap.~1).

\subsection*{Armadilhas conceituais frequentes}

\begin{itemize}
  \item \textbf{Sinais no HS}: $\zeta$ ciclônica é NEGATIVA (mesmo
        sinal de $f$); giro horário. Fixar a convenção cedo e desenhar
        tudo com o sul para baixo — importar figuras do HN sem espelhar
        é a maior fonte de confusão.
  \item \textbf{``A baixa suga o ar''}: a causalidade é da altitude
        para baixo — divergência em cima exporta massa, a pressão cai,
        e a convergência de superfície é a \emph{resposta}.
  \item \textbf{Ciclone = frente}: a frente é a estrutura; o ciclone é
        o modo instável que a usa. Eady não precisa de frente pronta.
  \item \textbf{Confundir advecção de vorticidade com vorticidade}: o
        ciclone de superfície fica sob o máximo de \emph{advecção} (a
        leste do cavado), não sob o cavado.
  \item \textbf{Esquecer o calor latente}: bombas são úmidas; a chuva
        frontal aquece a coluna e acelera a queda de pressão (T3).
\end{itemize}

\subsection*{Demonstração de baixo custo}

Balde girante com dreno central: abrir o dreno (convergência) e ver o
giro se intensificar — esticamento ao vivo. Cartas reais: pegar a
última ciclogênese do Atlântico Sul (GFS/GDAS em 500\,hPa +
superfície) e apontar cavado, advecção, streak e o ciclone na
superfície um quarto de onda a leste.

\subsection*{Problemas sugeridos do livro (exercícios do Cap.~13)}

Aquecimento: identificação de estágios do ciclo de vida em cartas;
esticamento com números dados. Aprofundamento: tendência de pressão por
camadas; quadrantes de streak para casos dados. Síntese: ``por que não
há ciclones extratropicais nos trópicos?'' (sem baroclinia nem $f$
suficiente). \emph{Confira a numeração na sua cópia (v1.02b).}

\subsection*{Uso do notebook \texttt{cap13\_ciclones.ipynb}}

\begin{center}
\begin{tabular}{lll}
\hline
Caso & Conteúdo & Momento sugerido \\
\hline
1 & $\zeta$ e advecção numa onda de 500\,hPa & Aula 1, após vorticidade \\
2 & taxa de Eady e $\lambda_{max}$ & Aula 1, fecho \\
3 & EDO de spin-up com/sem Ekman & Aula 2, abertura \\
4 & quadrantes do jet streak & Aula 2, após streaks \\
\hline
\end{tabular}
\end{center}

O Caso~3 é a ``previsão numérica em miniatura'' do capítulo: uma EDO
não linear que captura nascimento, teto e morte do ciclone.

\subsection*{Exercícios complementares e soluções}

Nas notas: T1--T5 e N1--N4. Soluções em \texttt{cap13\_solucoes.ipynb}
(\emph{não distribuir}): T1 verificado com \texttt{sympy}; N2 prevê a
temporada de ciclones do Atlântico Sul; N3 dá o ciclo de vida completo.
Pares: T1$\leftrightarrow$N1, T2$\leftrightarrow$N2,
T4$\leftrightarrow$N4.

\subsection*{Conexões com o resto do curso}

Transporte de energia (Cap.~11); frentes como estrutura (Cap.~12);
levantamento e nuvens estratiformes (Caps.~6, 10); calor latente
(Caps.~4, 7); o setor quente instável dispara os Caps.~14--15; a
previsão dessas ondas é o teste clássico dos modelos (Cap.~20).

\clearpage
% Guia do professor — Capítulo 14: Fundamentos de Tempestades
\section{Capítulo 14 --- Fundamentos de Tempestades}

\subsection*{Visão geral e ritmo}

O capítulo em que a termodinâmica dos Caps.~3--5 vira drama: CAPE, CIN,
gatilhos e o papel organizador do cisalhamento. A narrativa dos
``quatro ingredientes'' segura as duas aulas; o clímax é entender por
que o hodógrafo — e não o CAPE — decide o tipo de tempestade.
\textbf{Ritmo: 2 aulas de 2\,h.}

\begin{itemize}
  \item \textbf{Aula 1} — receita de 4 ingredientes; Skew-$T$ com CAPE
        e CIN (Caso~1 ao vivo); $w = \sqrt{2CAPE}$ e o fator
        $\sim$1/2 real; a tarde que estoura a tampa (N1).
  \item \textbf{Aula 2} — hodógrafos e os três tipos (Caso~3);
        supercélula em corte; piscina fria, frente de rajada e RKW
        (Casos~4); ciclo de vida e MCCs sul-americanos.
\end{itemize}

\subsection*{Objetivos de aprendizagem}

\begin{enumerate}
  \item calcular CAPE/CIN/LCL/LFC/EL de uma sondagem (MetPy e à mão no
        Skew-$T$);
  \item converter CAPE em $w_{max}$ e justificar o fator $\sim$1/2
        observado;
  \item classificar o regime convectivo pelo cisalhamento 0--6\,km e
        pelo BRN;
  \item explicar a auto-propagação via piscina fria (corrente de
        densidade) e o equilíbrio RKW;
  \item reconhecer as partes da supercélula (updraft, FFD/RFD,
        bigorna, overshooting top).
\end{enumerate}

\subsection*{Pré-requisitos a reativar}

Skew-$T$ e método da parcela (Cap.~5, essencial!); $T_v$ com carga de
água (Cap.~1); bulbo úmido e evaporação (Cap.~4); Bowen e o
combustível da tarde (Cap.~3); gatilhos frontais (Cap.~12).

\subsection*{Armadilhas conceituais frequentes}

\begin{itemize}
  \item \textbf{``CAPE grande $=$ tempestade severa''}: sem
        cisalhamento, CAPE de 3000 vira um pulso de 30\,min. Martelar
        o papel do hodógrafo.
  \item \textbf{CIN como vilão}: o CIN moderado é o que PERMITE a
        explosão da tarde — sem tampa não há panela de pressão.
  \item \textbf{Sentidos no HS}: hodógrafo útil curva ANTI-horário com
        a altura; mesociclones giram horário. Espelhar toda figura
        importada do HN.
  \item \textbf{$w = \sqrt{2CAPE}$ como previsão}: é teto teórico;
        entranhamento e água cobram metade (N2 mostra).
  \item \textbf{Confundir frente de rajada com frente fria}: mesma
        física de corrente de densidade (Cap.~12!), escalas e vidas
        diferentes.
\end{itemize}

\subsection*{Demonstração de baixo custo}

Corrente de densidade caseira: água salgada tingida despejada num
aquário de água doce — a ``frente de rajada'' com o nariz elevado
aparece perfeita. Vídeos de timelapse de supercélula (rotação visível
do updraft) valem 10 minutos de aula.

\subsection*{Problemas sugeridos do livro (exercícios do Cap.~14)}

Aquecimento: CAPE/CIN de sondagens dadas; $w_{max}$ e comparação com
updrafts típicos. Aprofundamento: BRN e classificação; velocidade de
frente de rajada. Síntese: ``por que o fim de tarde de verão é a hora
das tempestades no interior do Brasil?''.
\emph{Confira a numeração na sua cópia (v1.02b).}

\subsection*{Uso do notebook \texttt{cap14\_tempestades.ipynb}}

\begin{center}
\begin{tabular}{lll}
\hline
Caso & Conteúdo & Momento sugerido \\
\hline
1 & CAPE/CIN com MetPy (Skew-$T$ sombreado) & Aula 1, após parcela \\
2 & $w_{max}(\mathrm{CAPE})$ e o fator 1/2 & Aula 1 \\
3 & hodógrafos: única × multi × super & Aula 2, abertura \\
4 & piscina fria: $c(\Delta\theta_v, H)$ + downdraft & Aula 2 \\
\hline
\end{tabular}
\end{center}

O Caso~1 é o gancho perfeito para rodar ao vivo com a sondagem do dia
(Wyoming, SBMT 12\,UTC) e decidir com a turma ``vai ter tempestade
hoje?''.

\subsection*{Exercícios complementares e soluções}

Nas notas: T1--T5 e N1--N4. Soluções em \texttt{cap14\_solucoes.ipynb}
(\emph{não distribuir}): T2 com \texttt{sympy} (CAPE do triângulo);
N3 implementa SRH com movimento de tempestade do HS. Pares:
T1$\leftrightarrow$N2, T4$\leftrightarrow$N4, T5$\leftrightarrow$N3.

\subsection*{Conexões com o resto do curso}

Método da parcela e estabilidade (Cap.~5); microfísica da chuva e do
granizo (Cap.~7); assinaturas de radar (Cap.~8); gatilhos frontais e
linha seca (Cap.~12); os perigos (Cap.~15); convecção tropical e
furacões (Cap.~16); parametrização de convecção nos modelos (Cap.~20).

\clearpage
% Guia do professor — Capítulo 15: Perigos de Tempestades
\section{Capítulo 15 --- Perigos de Tempestades}

\subsection*{Visão geral e ritmo}

O capítulo mais ``vendável'' do curso — e uma coleção de física
clássica no extremo: balanço peso--arrasto (granizo), termodinâmica da
evaporação (downburst), eletrostática (raios) e ciclostrofia
(tornado). Cada perigo fecha com números de segurança reais.
\textbf{Ritmo: 2 aulas de 2\,h.}

\begin{itemize}
  \item \textbf{Aula 1} — granizo (balanço no quadro, biografia da
        pedra, derretimento na queda) e downbursts (DCAPE, seco ×
        úmido, microburst e aviação). Casos~1--2.
  \item \textbf{Aula 2} — relâmpagos (o capacitor no quadro,
        flash-to-bang, 30--30) e tornados (Rankine, queda de pressão,
        escala EF, lado perigoso no HS). Casos~3--4.
\end{itemize}

\subsection*{Objetivos de aprendizagem}

\begin{enumerate}
  \item calcular o updraft mínimo para cada tamanho de granizo;
  \item definir DCAPE e prever a intensidade de downbursts a partir da
        umidade sub-nuvem;
  \item descrever a eletrificação por colisões e usar as regras de
        segurança (3\,s/km, 30--30);
  \item integrar o equilíbrio ciclostrófico do Rankine e obter
        $\Delta P = \rho v_{max}^2$;
  \item relacionar a escala EF com a pressão dinâmica
        $\tfrac12\rho v^2$.
\end{enumerate}

\subsection*{Pré-requisitos a reativar}

Updraft e CAPE (Cap.~14); riming e fase mista (Cap.~7); $\theta_w$ e
evaporação (Cap.~4); equilíbrio ciclostrófico (Cap.~10); esticamento
de vorticidade (Cap.~13); refletividade e dual-pol (Cap.~8).

\subsection*{Armadilhas conceituais frequentes}

\begin{itemize}
  \item \textbf{``Granizo cai porque ficou pesado''}: cai quando
        $w_T(D)$ VENCE o updraft — a mesma pedra estaria voando numa
        supercélula mais forte.
  \item \textbf{Downburst exige chuva forte}: o seco cai de virga com
        eco fraco — o mais perigoso é o que não se vê.
  \item \textbf{Raio ``procura o mais alto''}: o líder desce por
        passos de campo máximo; altura ajuda mas não decide sozinha —
        daí a regra do abrigo, não da estatística.
  \item \textbf{Tornado como ``mini-ciclone''}: sem Coriolis! É
        ciclostrófico e pode girar para qualquer lado; o giro
        predominante vem do mesociclone-mãe.
  \item \textbf{Lado perigoso}: no HS é o ESQUERDO da trajetória
        (giro horário $+$ translação) — espelho dos livros do HN.
\end{itemize}

\subsection*{Demonstração de baixo custo}

Gerador eletrostático (Wimshurst/van de Graaff) para faíscas e o
conceito de ruptura; vórtice de garrafa PET (dois litros + água) para
o Rankine com ``olho'' visível. Vídeos: microburst de Andrews AFB
(anemômetro de 67\,m/s) e granizo em câmera lenta.

\subsection*{Problemas sugeridos do livro (exercícios do Cap.~15)}

Aquecimento: $w_T$ para tamanhos dados; flash-to-bang de tempos
medidos. Aprofundamento: DCAPE de sondagens; $\Delta P$ de tornados
com $v_{max}$ dados. Síntese: ``por que o corredor sul do Brasil /
norte da Argentina tem os MCCs mais elétricos do planeta?''.
\emph{Confira a numeração na sua cópia (v1.02b).}

\subsection*{Uso do notebook \texttt{cap15\_perigos.ipynb}}

\begin{center}
\begin{tabular}{lll}
\hline
Caso & Conteúdo & Momento sugerido \\
\hline
1 & $w_T(D)$ do granizo com marcos (ervilha$\to$laranja) & Aula 1 \\
2 & downburst: queda com déficit evaporativo & Aula 1 \\
3 & capacitor da tempestade: dente de serra de raios & Aula 2 \\
4 & Rankine: $v(r)$ e $\Delta P$ por EF & Aula 2 \\
\hline
\end{tabular}
\end{center}

\subsection*{Exercícios complementares e soluções}

Nas notas: T1--T5 e N1--N4. Soluções em \texttt{cap15\_solucoes.ipynb}
(\emph{não distribuir}): T1 e T4 com \texttt{sympy}; N1 e N4 são os
mais ricos (derretimento na queda; lado perigoso no HS). Pares:
T1$\leftrightarrow$N1, T2$\leftrightarrow$N2, T4$\leftrightarrow$N4.

\subsection*{Conexões com o resto do curso}

Fábrica (Cap.~14); microfísica do gelo (Cap.~7); radar e dual-pol
(Cap.~8); ciclostrofia (Cap.~10); esticamento (Cap.~13); ventos fortes
não-convectivos para contraste (Cap.~17); climatologia de severos num
clima mais quente (Cap.~21).

\clearpage
% Guia do professor — Capítulo 16: Ciclones Tropicais
\section{Capítulo 16 --- Ciclones Tropicais}

\subsection*{Visão geral e ritmo}

O anti-Cap.~13: núcleo quente, sem frentes, movido a calor latente do
oceano. A joia pedagógica é o furacão como máquina de Carnot — um
motor térmico desenhado no céu que a turma de física reconhece na
hora. Fechar com Catarina (2004) traz o tema para casa.
\textbf{Ritmo: 2 aulas de 2\,h.}

\begin{itemize}
  \item \textbf{Aula 1} — anatomia (olho, parede, bandas), tabela de
        contraste com o Cap.~13, condições de gênese, Carnot/MPI e
        WISHE. Casos~1 e 3 (o \texttt{tcpyPI} impressiona).
  \item \textbf{Aula 2} — Holland e vento gradiente no limite,
        trajetórias e deriva beta, perigos com ênfase em surge
        (Caso~4), Catarina e o clima futuro (N3).
\end{itemize}

\subsection*{Objetivos de aprendizagem}

\begin{enumerate}
  \item contrastar ciclones tropicais e extratropicais (energia,
        núcleo, estrutura);
  \item enunciar as condições de gênese e usá-las para explicar a
        climatologia das bacias;
  \item calcular a intensidade potencial (Carnot simplificado e
        \texttt{tcpyPI}) e explicar o WISHE;
  \item obter o perfil de vento do modelo de Holland;
  \item estimar maré de tempestade (empilhamento $+$ barômetro
        invertido) e ranquear os perigos.
\end{enumerate}

\subsection*{Pré-requisitos a reativar}

Clausius--Clapeyron e $\theta_e$ (Cap.~4); fluxos bulk (Cap.~3);
vento gradiente e Coriolis (Cap.~10); hipsométrica para o núcleo
quente (Cap.~1); alísios e $\beta$ (Cap.~11); CAPE (Cap.~14).

\subsection*{Armadilhas conceituais frequentes}

\begin{itemize}
  \item \textbf{``Furacão é tempestade grande''}: é um sistema
        AXISSIMÉTRICO de núcleo quente — outra classe de objeto; as
        tempestades do Cap.~14 são os tijolos da parede.
  \item \textbf{Cisalhamento}: no Cap.~14 organiza; aqui MATA
        (ventila o núcleo quente). Vale uma linha no quadro.
  \item \textbf{``26,5\,°C é lei''}: é limiar empírico do $\Delta k$;
        com topo muito frio, menos serve (subtropicais, medicanes,
        Catarina).
  \item \textbf{Olho calmo $=$ fim}: a parede oposta chega em
        minutos — erro clássico (e fatal) de quem sai do abrigo.
  \item \textbf{Vento como maior perigo}: a água (surge $+$ chuva)
        mata mais; dano eólico $\propto v^3$, mas enchente não segue
        categoria.
\end{itemize}

\subsection*{Demonstração de baixo custo}

Ciclo de Carnot desenhado sobre o corte do furacão (transparência ou
quadro em duas cores) — a analogia motor térmico rende. Imagens de
satélite do Catarina (2004) e discussão: por que surpreendeu?
(Atlântico Sul ``não tinha'' furacões.)

\subsection*{Problemas sugeridos do livro (exercícios do Cap.~16)}

Aquecimento: MPI para TSMs dadas; queda de pressão de núcleo quente.
Aprofundamento: vento gradiente de perfis dados; surge com plataforma
dada. Síntese: ``por que não há furacões no Atlântico Sul — e o que o
Catarina ensinou?''. \emph{Confira a numeração na sua cópia (v1.02b).}

\subsection*{Uso do notebook \texttt{cap16\_furacoes.ipynb}}

\begin{center}
\begin{tabular}{lll}
\hline
Caso & Conteúdo & Momento sugerido \\
\hline
1 & Carnot + Clausius--Clapeyron: MPI(TSM) & Aula 1 \\
2 & Holland: $P(r)$ e $v(r)$ por categoria & Aula 2 \\
3 & \texttt{tcpyPI}: MPI de sondagem real & Aula 1, fecho \\
4 & surge: plataforma rasa × profunda & Aula 2 \\
\hline
\end{tabular}
\end{center}

O Caso~3 usa o pacote de pesquisa real (Bister--Emanuel) — vale
mostrar que a ferramenta de sala é a mesma dos artigos.

\subsection*{Exercícios complementares e soluções}

Nas notas: T1--T5 e N1--N4. Soluções em \texttt{cap16\_solucoes.ipynb}
(\emph{não distribuir}): T2 com \texttt{sympy}; N2 ajusta Holland por
otimização; N3 é a ponte quantitativa para o Cap.~21. Pares:
T1$\leftrightarrow$N1, T2$\leftrightarrow$N2, T4$\leftrightarrow$N4.

\subsection*{Conexões com o resto do curso}

Clausius--Clapeyron e calor latente (Cap.~4); fluxos ar--mar
(Cap.~3); vento gradiente (Cap.~10); alísios, $\beta$ e recurvamento
(Cap.~11); convecção da parede (Cap.~14); MPI num clima quente
(Cap.~21); previsão de trajetória por ensemble (Cap.~20).

\clearpage
% Guia do professor — Capítulo 17: Ventos Regionais
\section{Capítulo 17 --- Ventos Regionais}

\subsection*{Visão geral e ritmo}

O capítulo da geografia local: brisas, vales, catabáticos, ondas de
montanha e a hidráulica de sotavento. Fisicamente delicioso — teorema
de circulação, oscilador $N$, água rasa e Froude — e imediatamente
reconhecível pela turma (Santos, serra do Mar, vale do Paraíba).
\textbf{Ritmo: 2 aulas de 2\,h.}

\begin{itemize}
  \item \textbf{Aula 1} — Bjerknes e a brisa (com o hodógrafo do
        Caso~1), terral, ventos de vale/encosta, catabáticos até a
        Antártida (Caso~2).
  \item \textbf{Aula 2} — ondas de montanha e Scorer (Caso~3), föhn
        completo, água rasa/Froude e vendavais (Caso~4), gap winds.
\end{itemize}

\subsection*{Objetivos de aprendizagem}

\begin{enumerate}
  \item explicar a brisa pelo teorema de circulação e prever seu giro
        diurno no HS;
  \item calcular a velocidade de equilíbrio catabática;
  \item usar o parâmetro de Scorer para classificar ondas
        (propagante/evanescente/aprisionada) e obter $\lambda_z$;
  \item calcular o aquecimento föhn com remoção de chuva;
  \item usar o número de Froude para prever bloqueio, travessia e
        vendaval de sotavento.
\end{enumerate}

\subsection*{Pré-requisitos a reativar}

Bowen e contraste terra--mar (Cap.~3); resfriamento radiativo noturno
(Caps.~2, 6); $N$ e estabilidade (Cap.~5); $\theta$ e föhn (Cap.~3);
nuvens de onda (Cap.~6); Coriolis (Cap.~10).

\subsection*{Armadilhas conceituais frequentes}

\begin{itemize}
  \item \textbf{``Brisa é o ar indo do frio pro quente''}: embaixo
        sim, mas é um CIRCUITO — sem o ramo de retorno em altitude
        não fecha; o teorema de Bjerknes resolve.
  \item \textbf{Giro da brisa}: no HS o vetor gira anti-horário ao
        longo do dia (espelho do HN).
  \item \textbf{Lenticular como nuvem que anda}: está PARADA sobre a
        crista com vento forte através — ótimo choque conceitual.
  \item \textbf{Föhn como ``compressão''}: o ar não é comprimido como
        pistão; desce pelo $\Gamma_d$ porque a chuva removeu água e
        $L_v$ a barlavento (T5 fecha a conta).
  \item \textbf{Vendaval de sotavento como ``vento que despenca''}: é
        transição hidráulica (Fr), não queda livre — a analogia da
        pia é a ponte certa.
\end{itemize}

\subsection*{Demonstração de baixo custo}

A pia da cozinha: jato supercrítico + salto hidráulico em anel —
todos já viram, ninguém tinha nomeado. Tanque com água e obstáculo
(régua) puxado a velocidades diferentes: os três regimes de Froude
aparecem. Vídeos: lenticulares em timelapse; Bora em Trieste.

\subsection*{Problemas sugeridos do livro (exercícios do Cap.~17)}

Aquecimento: $u_{eq}$ catabático com números dados; $\lambda_z$ de
ondas. Aprofundamento: föhn com LCL dado; Froude de camadas dadas.
Síntese: ``desenhe o dia de ventos de uma cidade costeira com serra
atrás (Santos!): brisa, terral, catabático e onda no mesmo ciclo''.
\emph{Confira a numeração na sua cópia (v1.02b).}

\subsection*{Uso do notebook \texttt{cap17\_ventosregionais.ipynb}}

\begin{center}
\begin{tabular}{lll}
\hline
Caso & Conteúdo & Momento sugerido \\
\hline
1 & brisa: EDO com Coriolis + hodógrafo diurno & Aula 1 \\
2 & catabático: EDO e $u_{eq}(\alpha, \Delta\theta)$ & Aula 1 \\
3 & ondas de montanha lineares (morro-sino) & Aula 2 \\
4 & água rasa sobre a serra: regimes de Froude & Aula 2 \\
\hline
\end{tabular}
\end{center}

\subsection*{Exercícios complementares e soluções}

Nas notas: T1--T5 e N1--N4. Soluções em \texttt{cap17\_solucoes.ipynb}
(\emph{não distribuir}): T2 com \texttt{sympy}; N1 responde quando a
brisa perde do vento sinótico; N4 constrói o mapa de regimes. Pares:
T2$\leftrightarrow$N2, T3$\leftrightarrow$N3, T4$\leftrightarrow$N4.

\subsection*{Conexões com o resto do curso}

Circuitos térmicos e Bowen (Cap.~3); $N$ como mola (Cap.~5);
lenticulares e rotores (Cap.~6); brisa como gatilho de convecção
(Cap.~14); catabático e piscina fria noturna (Cap.~18); arrasto de
onda nos modelos (Cap.~20); brisas e ilha de calor urbana (Cap.~19).

\clearpage
% Guia do professor — Capítulo 18: Camada Limite Atmosférica
\section{Capítulo 18 --- Camada Limite Atmosférica}

\subsection*{Visão geral e ritmo}

O capítulo em que a turma descobre que mora dentro do objeto de
estudo. Três argumentos dimensionais (perfil log, orçamento de TKE,
$-5/3$) fazem a ponte perfeita com a formação em física; o ciclo
diurno costura tudo. \textbf{Ritmo: 2 aulas de 2\,h.}

\begin{itemize}
  \item \textbf{Aula 1} — o ``dia da CLA'' (desenho-mestre),
        encroachment e $z_i$ (Caso~1), camada superficial e o log
        (Caso~2, com energia eólica).
  \item \textbf{Aula 2} — TKE e seus termos, Kolmogorov com o
        espectro sintético (Caso~3), a noite estável e o jato de
        Blackadar (Caso~4).
\end{itemize}

\subsection*{Objetivos de aprendizagem}

\begin{enumerate}
  \item narrar o ciclo diurno da CLA (mistura, residual, estável) e
        seus perfis de $\theta_v$;
  \item prever $z_i(t)$ pelo encroachment a partir de $F_H$ e
        $\gamma$;
  \item ajustar o perfil log ($u_*$, $z_0$) a dados de torre e
        aplicá-lo a energia eólica;
  \item interpretar o orçamento de TKE e o espectro $-5/3$;
  \item explicar o jato noturno como oscilação inercial de Blackadar.
\end{enumerate}

\subsection*{Pré-requisitos a reativar}

Balanço de superfície e Bowen (Cap.~3); perfis de $\theta_v$ e
estabilidade não local (Cap.~5); círculos inerciais e Ekman
(Cap.~10); resfriamento radiativo noturno (Caps.~2, 6); $Ri$
(Cap.~5).

\subsection*{Armadilhas conceituais frequentes}

\begin{itemize}
  \item \textbf{``A CLA é fina, logo irrelevante''}: é onde vive TODO
        o vapor dos Caps.~4--7, o combustível do Cap.~14 e a poluição
        do Cap.~19 — e onde o vento vira recurso.
  \item \textbf{Perfil log como lei universal}: vale na camada
        superficial neutra; estabilidade curva o perfil
        (Monin--Obukhov). Não extrapolar o log até $z_i$.
  \item \textbf{Turbulência como ruído}: tem orçamento, cascata e
        espectro — é física estatística organizada, não sujeira.
  \item \textbf{Jato noturno por ``canalização''}: é inércia pura
        (Blackadar); o hodógrafo circular do Caso~4 desfaz o mito.
  \item \textbf{Camada residual esquecida}: o ar da tarde de ontem
        sobrevive à noite em cima — e é o que a térmica de amanhã
        reencontra (e o que guarda poluição envelhecida, Cap.~19).
\end{itemize}

\subsection*{Demonstração de baixo custo}

Timelapse de cumulus de bom tempo (o teto plano dos cumulus $=$
$z_i$); chaleira fumegante ao sol da janela para térmicas;
anemômetro de celular/ventoinha em alturas diferentes no pátio para o
gradiente do perfil (qualitativo, mas memorável).

\subsection*{Problemas sugeridos do livro (exercícios do Cap.~18)}

Aquecimento: $z_i$ por encroachment com números dados; $u(z)$ do log.
Aprofundamento: TKE com fluxos dados; jato noturno com $f$ de
latitudes diferentes. Síntese: ``por que planadores decolam às 11\,h
e pousam às 17\,h?''. \emph{Confira a numeração na sua cópia
(v1.02b).}

\subsection*{Uso do notebook \texttt{cap18\_camadalimite.ipynb}}

\begin{center}
\begin{tabular}{lll}
\hline
Caso & Conteúdo & Momento sugerido \\
\hline
1 & crescimento de $z_i$ (EDO vs.\ raiz) & Aula 1 \\
2 & ajuste do perfil log $+$ potência eólica & Aula 1 \\
3 & espectro $-5/3$ de sinal sintético (FFT) & Aula 2 \\
4 & jato de Blackadar (hodógrafo circular) & Aula 2 \\
\hline
\end{tabular}
\end{center}

O Caso~3 rende discussão de método experimental (taxa de amostragem,
aliasing — N3): meteorologia como física de laboratório.

\subsection*{Exercícios complementares e soluções}

Nas notas: T1--T5 e N1--N4. Soluções em \texttt{cap18\_solucoes.ipynb}
(\emph{não distribuir}): T1 e T4 com \texttt{sympy}; N4 explica por
que o LLJ é fenômeno subtropical/de latitudes médias. Pares:
T1$\leftrightarrow$N1, T2$\leftrightarrow$N2, T4$\leftrightarrow$N4,
T5$\leftrightarrow$N3.

\subsection*{Conexões com o resto do curso}

Balanço de superfície (Cap.~3); perfis e estabilidade (Cap.~5);
entranhamento (Cap.~14); Ekman e spin-down (Caps.~10, 13); brisas e
catabáticos como CLAs especiais (Cap.~17); dispersão de poluentes usa
TUDO deste capítulo (Cap.~19); parametrização de CLA nos modelos
(Cap.~20).

\clearpage
% Guia do professor — Capítulo 19: Dispersão de Poluentes
\section{Capítulo 19 --- Dispersão de Poluentes}

\subsection*{Visão geral e ritmo}

O Cap.~18 aplicado a um problema com consequência social imediata.
Para físicos é um banquete de métodos: equação do calor, método das
imagens, balanço de caixa e movimento browniano — cada um resolvendo
um pedaço do problema regulatório. \textbf{Ritmo: 2 aulas de 2\,h
(cabe em 1 aula $+$ estudo dirigido se o semestre apertar).}

\begin{itemize}
  \item \textbf{Aula 1} — pluma gaussiana (dedução + imagens),
        máximo no solo e o $1/H^2$, as três plumas da estabilidade
        (Casos~1--2).
  \item \textbf{Aula 2} — modelo de caixa e episódios (Caso~3),
        fumigação, passeio aleatório e Taylor (Caso~4).
\end{itemize}

\subsection*{Objetivos de aprendizagem}

\begin{enumerate}
  \item deduzir a pluma gaussiana e aplicar a fonte-imagem;
  \item localizar e escalar o máximo de concentração no solo;
  \item associar forma de pluma a estabilidade e hora do dia;
  \item usar o modelo de caixa e o fator de ventilação para explicar
        episódios;
  \item conectar passeio aleatório, teorema de Taylor e difusividade
        turbulenta.
\end{enumerate}

\subsection*{Pré-requisitos a reativar}

CLA completa (Cap.~18: $z_i$, estabilidade, TKE, $T_L$); perfis e
inversões (Cap.~5); método das imagens (eletrostática!); equação do
calor; ventilação anticiclônica (Cap.~12).

\subsection*{Armadilhas conceituais frequentes}

\begin{itemize}
  \item \textbf{``Chaminé alta resolve''}: reduz o pico local
        ($1/H^2$) e EXPORTA o problema — gancho histórico da chuva
        ácida.
  \item \textbf{``Episódio $=$ mais emissão''}: os episódios de
        inverno são queda do fator de ventilação com emissão igual —
        o Caso~3 desfaz o mito.
  \item \textbf{Pluma estável ``comportada''}: o fanning é uma bomba
        armada esperando a fumigação da manhã.
  \item \textbf{Gaussiana como lei exata}: vale para terreno plano,
        estacionário, sem química; é o LEGO básico dos modelos, não
        o edifício.
  \item \textbf{Máximo noturno de poluição atribuído ao tráfego}: o
        rush é às 18\,h; o pico de C é mais tarde porque a caixa
        encolhe (N3).
\end{itemize}

\subsection*{Demonstração de baixo custo}

Incenso em sala: perto do teto (estável, estratificado) vira fanning;
sobre um aquecedor, looping — as três plumas por R\$\,5. Dados
CETESB do último inverno: sobrepor série de MP$_{10}$ com $U$ e
discutir o fator de ventilação.

\subsection*{Problemas sugeridos do livro (exercícios do Cap.~19)}

Aquecimento: $C$ no solo para fontes dadas; classes de Pasquill.
Aprofundamento: $x_{max}$, caixa com números reais. Síntese: ``por
que os piores episódios de SP são em julho, se se dirige o ano
todo?''. \emph{Confira a numeração na sua cópia (v1.02b).}

\subsection*{Uso do notebook \texttt{cap19\_poluentes.ipynb}}

\begin{center}
\begin{tabular}{lll}
\hline
Caso & Conteúdo & Momento sugerido \\
\hline
1 & pluma gaussiana: campo e pico no solo & Aula 1 \\
2 & looping/coning/fanning & Aula 1 \\
3 & caixa da cidade: 3 ventilações & Aula 2 \\
4 & Langevin: balístico $\to$ difusivo & Aula 2 \\
\hline
\end{tabular}
\end{center}

O Caso~4 fecha o arco ``física estatística na atmosfera'' começado no
Cap.~18 — Einstein 1905 com turbulência no lugar de moléculas.

\subsection*{Exercícios complementares e soluções}

Nas notas: T1--T5 e N1--N4. Soluções em \texttt{cap19\_solucoes.ipynb}
(\emph{não distribuir}): T3 com \texttt{sympy}; N2 (fumigação) e N3
(máximo noturno) são os de maior rendimento conceitual. Pares:
T3$\leftrightarrow$N1, T4$\leftrightarrow$N3, T5$\leftrightarrow$N4.

\subsection*{Conexões com o resto do curso}

CLA e turbulência (Cap.~18, essencial); estabilidade (Cap.~5);
inversões e anticiclones (Caps.~5, 12); brisas redistribuindo poluição
(Cap.~17); deposição e nuvens (Caps.~6--7); química do ozônio e clima
(Cap.~21); os mesmos esquemas de advecção nos modelos (Cap.~20).

\clearpage
% Guia do professor — Capítulo 20: Previsão Numérica do Tempo
\section{Capítulo 20 --- Previsão Numérica do Tempo}

\subsection*{Visão geral e ritmo}

O capítulo-síntese computacional: as 7 equações do curso viram
código, e as quatro peças (discretizar, parametrizar, inicializar,
quantificar) organizam tudo. Para físicos é território natal —
métodos numéricos, caos, inferência — aplicado ao maior problema de
valor inicial do mundo. \textbf{Ritmo: 2 aulas de 2\,h.}

\begin{itemize}
  \item \textbf{Aula 1} — Richardson$\to$ENIAC (10 min de história
        que valem ouro), grade, esquemas e CFL (Caso~1 ao vivo),
        parametrizações como física reduzida.
  \item \textbf{Aula 2} — Lorenz e o horizonte (Caso~2), assimilação
        (Caso~3), ensembles e probabilidade (Caso~4), a revolução
        silenciosa.
\end{itemize}

\subsection*{Objetivos de aprendizagem}

\begin{enumerate}
  \item listar as equações primitivas e reconhecê-las como o curso
        inteiro;
  \item aplicar CFL e prever os erros característicos de cada
        esquema;
  \item explicar o horizonte de ~2 semanas via Lyapunov;
  \item usar o ganho ótimo da assimilação e explicar por que a
        análise supera modelo e observação;
  \item interpretar previsões probabilísticas de ensemble.
\end{enumerate}

\subsection*{Pré-requisitos a reativar}

As 7 equações espalhadas (Caps.~1, 3, 4, 10); advecção numérica
(Cap.~3, N3!); tudo que virou parametrização (Caps.~2, 7, 14, 17,
18); mínimos quadrados/propagação de erros (laboratório de física).

\subsection*{Armadilhas conceituais frequentes}

\begin{itemize}
  \item \textbf{``Mais computador $=$ mais previsão''}: o horizonte é
        logarítmico no erro inicial (T3) — caos é estrutural, não
        falta de hardware.
  \item \textbf{``O modelo de 10\,km resolve 10\,km''}: resolução
        efetiva $\sim7\Delta x$; cumulus continuam parametrizados.
  \item \textbf{Assimilação como ``usar os dados''}: é média
        ponderada por precisões — jogar o modelo fora e usar só
        observações seria PIOR (T4).
  \item \textbf{``70\% de chuva'' como enfeite}: é a fração literal
        do ensemble — probabilidade é a saída primária, não o
        determinístico.
  \item \textbf{Caos como ruído}: o atrator é determinístico e
        estruturado — imprevisível no ponto, previsível na
        estatística (ponte para o Cap.~21).
\end{itemize}

\subsection*{Demonstração de baixo custo}

Rodar o Caso~2 ao vivo mudando o 6º decimal — o momento
``uau'' garantido do semestre. Comparar a previsão de 5 dias do
ECMWF/GFS de ontem com a carta de hoje; abrir o meteograma de
ensemble do ECMWF para a cidade (espaguete real).

\subsection*{Problemas sugeridos do livro (exercícios do Cap.~20)}

Aquecimento: CFL para grades/ventos dados; contas de pontos de grade
e memória. Aprofundamento: von Neumann de outros esquemas; ganho
ótimo com números. Síntese: ``por que o clima de 2100 é mais
previsível que o tempo de daqui a 20 dias?''.
\emph{Confira a numeração na sua cópia (v1.02b).}

\subsection*{Uso do notebook \texttt{cap20\_previsao.ipynb}}

\begin{center}
\begin{tabular}{lll}
\hline
Caso & Conteúdo & Momento sugerido \\
\hline
1 & upwind/Lax--Wendroff/leapfrog $+$ CFL & Aula 1 \\
2 & Lorenz 63: borboleta e gêmeos & Aula 2, abertura \\
3 & ciclo de assimilação no Lorenz & Aula 2 \\
4 & ensemble: espaguete $\to$ probabilidade & Aula 2, fecho \\
\hline
\end{tabular}
\end{center}

O arco dos Casos 2$\to$3$\to$4 é a história intelectual da previsão
moderna em 30 minutos de aula.

\subsection*{Exercícios complementares e soluções}

Nas notas: T1--T5 e N1--N4. Soluções em \texttt{cap20\_solucoes.ipynb}
(\emph{não distribuir}): T1, T2 e T4 com \texttt{sympy} (von Neumann,
equação modificada e ganho ótimo simbólicos); N4 é a pérola
(dispersão prevê o erro). Pares: T1/T2$\leftrightarrow$N1,
T3$\leftrightarrow$N2, T4$\leftrightarrow$N3.

\subsection*{Conexões com o resto do curso}

TODAS: as equações (Caps.~1, 3, 4, 10), as parametrizações
(Caps.~2, 7, 14, 17, 18), as observações (Caps.~8--9), os fenômenos
que o modelo deve acertar (Caps.~11--16) e a fronteira com o clima
(Cap.~21). Este capítulo é o ponto de fuga do curso.

\clearpage
% Guia do professor — Capítulo 21: Processos Climáticos Naturais
\section{Capítulo 21 --- Processos Climáticos Naturais}

\subsection*{Visão geral e ritmo}

A ponte tempo$\to$clima: feedbacks (a álgebra do amplificador),
escalas lentas e variabilidade interna. Quatro modelos de brinquedo
carregam tudo — gelo-albedo, Milankovitch, Daisyworld e o oscilador
ENSO — e cada um é um clássico da literatura rodando no notebook.
\textbf{Ritmo: 2 aulas de 2\,h.}

\begin{itemize}
  \item \textbf{Aula 1} — tempo × clima (por que a estatística é
        previsível), álgebra de feedbacks, bifurcação gelo-albedo
        (Caso~1 — o gráfico da aula), Milankovitch (Caso~2).
  \item \textbf{Aula 2} — Daisyworld (Caso~3), variabilidade interna
        e ENSO (Caso~4), vulcões e Sol, fechamento para o módulo extra (Cap.~23).
\end{itemize}

\subsection*{Objetivos de aprendizagem}

\begin{enumerate}
  \item usar a álgebra de feedbacks ($\Delta T =
        \lambda_0\Delta F/(1-f)$) e classificar os feedbacks
        principais;
  \item explicar múltiplos equilíbrios e histerese do gelo-albedo;
  \item associar os três ciclos orbitais a seus períodos e ao gatilho
        de 65°N;
  \item interpretar o Daisyworld como feedback negativo emergente;
  \item descrever o mecanismo do ENSO e sua previsibilidade sazonal.
\end{enumerate}

\subsection*{Pré-requisitos a reativar}

Balanço radiativo e $T_e$ (Cap.~2); Clausius--Clapeyron para o
feedback do vapor (Cap.~4); insolação orbital (Cap.~2, Caso~2);
alísios e Walker (Cap.~11); caos vs.\ estatística (Cap.~20);
realimentação (eletrônica básica).

\subsection*{Armadilhas conceituais frequentes}

\begin{itemize}
  \item \textbf{``O clima sempre mudou, logo\ldots''}: sim — e cada
        mudança teve CAUSA física identificável (órbita, vulcões,
        CO$_2$); entender as naturais é o que permite atribuir a
        atual (Cap.~23).
  \item \textbf{Feedback positivo $=$ explosão}: com $f<1$ é
        amplificação finita; a fuga só ocorre em $f\to1$ (bola de
        neve).
  \item \textbf{Milankovitch como ``causa direta''}: é o metrônomo;
        a amplitude vem dos feedbacks (gelo, CO$_2$) — N2 mostra a
        retificação.
  \item \textbf{Daisyworld como Gaia teleológica}: não há intenção —
        só seleção $+$ acoplamento; o modelo existe exatamente para
        desfazer essa confusão.
  \item \textbf{ENSO como ciclo regular}: é oscilador irregular
        (ruído $+$ não linearidade) — 3--7 anos, não relógio.
\end{itemize}

\subsection*{Demonstração de baixo custo}

A bifurcação do Caso~1 desenhada no quadro com setas de estabilidade
(análise gráfica de EDO — os alunos fazem sozinhos). Vídeo do tanque
giratório de convecção ou da série de TSM do Niño-3.4 (ondas reais de
El Niño/La Niña na tela).

\subsection*{Problemas sugeridos do livro (exercícios do Cap.~21)}

Aquecimento: resposta de Planck com números; períodos orbitais.
Aprofundamento: ganho com múltiplos feedbacks; equilíbrios do EBM com
albedo degrau. Síntese: ``por que o clima é previsível se o tempo não
é?''. \emph{Confira a numeração na sua cópia (v1.02b).}

\subsection*{Uso do notebook \texttt{cap21\_clima.ipynb}}

\begin{center}
\begin{tabular}{lll}
\hline
Caso & Conteúdo & Momento sugerido \\
\hline
1 & bifurcação gelo-albedo (bola de neve) & Aula 1 \\
2 & ciclos de Milankovitch e $Q_{65N}$ & Aula 1 \\
3 & Daisyworld completo & Aula 2 \\
4 & oscilador de recarga do ENSO & Aula 2 \\
\hline
\end{tabular}
\end{center}

Os quatro casos são papers clássicos (Budyko/Sellers 1969; Milankovitch;
Watson--Lovelock 1983; Jin 1997) em miniatura — vale dizer isso à
turma: eles estão rodando a literatura.

\subsection*{Exercícios complementares e soluções}

Nas notas: T1--T5 e N1--N4. Soluções em \texttt{cap21\_solucoes.ipynb}
(\emph{não distribuir}): T2 e T5 com \texttt{sympy}; N1 e N2 são os
mais ricos (histerese; retificação de Milankovitch). Pares:
T1$\leftrightarrow$N1, T3$\leftrightarrow$N2, T4$\leftrightarrow$N3,
T5$\leftrightarrow$N4.

\subsection*{Conexões com o resto do curso}

Balanço radiativo (Cap.~2); vapor e Clausius--Clapeyron (Cap.~4);
nuvens como feedback incerto (Caps.~6--7); circulação de Walker e
alísios (Cap.~11); caos e previsibilidade estatística (Cap.~20); o
Cap.~23 (módulo extra) aplica TUDO isto à perturbação antrópica; o Cap.~22 fecha o livro com a óptica atmosférica.

\clearpage
% Guia do professor — Capítulo 22: Óptica Atmosférica
\section{Capítulo 22 --- Óptica Atmosférica}

\subsection*{Visão geral e ritmo}

O capítulo-sobremesa do livro: Snell, dispersão, dipolos e refração
aplicados ao céu. Para físicos é território natal — e a mensagem
docente é que cada fenômeno luminoso é um \emph{diagnóstico} (halo $=$
cirrus; arco $=$ gotas grandes; miragem $=$ perfil de $T$).
\textbf{Ritmo: 2 aulas de 2\,h (rende bem como fecho leve do
semestre).}

\begin{itemize}
  \item \textbf{Aula 1} — arco-íris completo (cáustica no quadro,
        Monte Carlo no Caso~1), halos e a família dos cristais
        (Caso~2).
  \item \textbf{Aula 2} — Rayleigh e as cores do céu (Caso~3),
        polarização, refração e miragens (Caso~4), flash verde, e a
        síntese do curso.
\end{itemize}

\subsection*{Objetivos de aprendizagem}

\begin{enumerate}
  \item deduzir o ângulo do arco-íris pela cáustica do desvio mínimo;
  \item explicar o halo de 22° pelo desvio mínimo do prisma de gelo;
  \item usar Rayleigh ($\lambda^{-4}$) $+$ Beer para céu azul e poente
        vermelho;
  \item calcular a curvatura de raios em gradientes de $n$ e
        classificar miragens;
  \item ler fenômenos ópticos como diagnósticos de nuvens e perfis.
\end{enumerate}

\subsection*{Pré-requisitos a reativar}

Snell e dispersão (óptica básica); Beer e massa de ar (Cap.~2); gotas
e cristais e seus tamanhos (Caps.~6--7); inversões (Cap.~5); lei dos
gases para $n(\rho)$ (Cap.~1); dipolo irradiante (eletromagnetismo).

\subsection*{Armadilhas conceituais frequentes}

\begin{itemize}
  \item \textbf{``O arco está num lugar''}: é um cone de 42° em torno
        do antissolar de CADA observador — não há pé do arco-íris.
  \item \textbf{Halo confundido com coroa}: halo $=$ refração em gelo
        (22°, borda vermelha interna); coroa $=$ difração em gotículas
        (poucos graus, vermelho externo).
  \item \textbf{``Céu azul porque reflete o mar''}: o dipolo de
        Rayleigh resolve em uma linha (T3) — e a polarização (T5) é a
        prova experimental barata (óculos polarizados!).
  \item \textbf{Miragem como ``ilusão''}: os raios chegam de verdade
        por caminhos curvos — a imagem é real, a interpretação do
        cérebro é que assume linha reta.
  \item \textbf{Flash verde como lenda}: é dispersão $+$ refração
        padrão; raro porque exige horizonte limpo e baixo aerossol.
\end{itemize}

\subsection*{Demonstração de baixo custo}

Laser $+$ aquário com gelatina/açúcar em gradiente: o raio curva à
vista (miragem de bancada). Óculos polarizados girando contra o céu a
90° do Sol. Esfera de vidro/gota pendurada ao sol para o arco de
bancada. CD como rede de difração para o espectro do céu.

\subsection*{Problemas sugeridos do livro (exercícios do Cap.~22)}

Aquecimento: ângulos de arco e halo com $n$ dados; massa de ar do
poente. Aprofundamento: curvatura de raios com gradientes dados;
polarização. Síntese: ``você vê um halo pela manhã — que nuvem é, o
que vem nas próximas 24\,h e por quê?'' (amarra Caps.~6, 12 e 22).
\emph{Confira a numeração na sua cópia (v1.02b).}

\subsection*{Uso do notebook \texttt{cap22\_optica.ipynb}}

\begin{center}
\begin{tabular}{lll}
\hline
Caso & Conteúdo & Momento sugerido \\
\hline
1 & arco-íris: traçado de raios $+$ Monte Carlo & Aula 1 \\
2 & halo 22°: prisma e histograma de cristais & Aula 1 \\
3 & Rayleigh $+$ Beer: céu e poente & Aula 2 \\
4 & miragens: raios integrados em $n(z)$ & Aula 2 \\
\hline
\end{tabular}
\end{center}

Os Casos 1--2 são Monte Carlo de verdade com física de graduação — um
fecho computacional elegante para o curso.

\subsection*{Exercícios complementares e soluções}

Nas notas: T1--T5 e N1--N4. Soluções em \texttt{cap22\_solucoes.ipynb}
(\emph{não distribuir}): T1 e T2 com \texttt{sympy}; N4 (looming) é o
mais bonito. Pares: T1$\leftrightarrow$N1, T2$\leftrightarrow$N2,
T3$\leftrightarrow$N3, T4$\leftrightarrow$N4.

\subsection*{Conexões com o resto do curso}

Beer e massa de ar (Cap.~2); gotas e cristais (Caps.~6--7: o arco
exige chuva, o halo exige cirrus); halo como precursor frontal
(Cap.~12); inversões e miragens (Cap.~5); turbulência e cintilação
(Cap.~18); aerossol vulcânico e poentes (Cap.~21). O módulo extra de
Mudança Climática (Cap.~23) pode fechar o semestre após este.

\clearpage
% Guia do professor — Capítulo 23 (módulo extra): Mudança Climática
\section{Capítulo 23 (módulo extra) --- Mudança Climática}

\subsection*{Visão geral e ritmo}

O fecho do curso: nada de novo entra — estufa, feedbacks, oceano e
estatística já foram construídos; aqui são aplicados à perturbação
antrópica com números na mesa. Tom recomendado: física serena e
quantitativa; deixar os dados falarem. \textbf{Ritmo: 2 aulas de 2\,h
(a última pode fechar com retrospectiva do curso).}

\begin{itemize}
  \item \textbf{Aula 1} — o logaritmo do CO$_2$ (T1 no quadro),
        sensibilidade e nuvens, a inércia das duas caixas (Casos~1--2).
  \item \textbf{Aula 2} — TCRE e orçamento (Caso~3), detecção no
        ruído (Caso~4), impressões digitais (estratosfera!), impactos
        capítulo a capítulo e o fecho do curso.
\end{itemize}

\subsection*{Objetivos de aprendizagem}

\begin{enumerate}
  \item deduzir/justificar $\Delta F = 5{,}35\ln(C/C_0)$ e calcular a
        forçante atual;
  \item combinar forçante e feedbacks para a faixa de sensibilidade;
  \item explicar o aquecimento comprometido pelas escalas oceânicas;
  \item usar o TCRE para converter metas em orçamentos de carbono;
  \item distinguir sinal de ruído e listar as assinaturas do
        mecanismo estufa.
\end{enumerate}

\subsection*{Pré-requisitos a reativar}

Bandas e janelas (Cap.~2); feedbacks e $\lambda_0$ (Cap.~21);
capacidade térmica do oceano (Cap.~3); ENSO como ruído (Cap.~21);
tendência vs.\ variância (estatística básica); MPI $+2$\,K (Cap.~16,
N3).

\subsection*{Armadilhas conceituais frequentes}

\begin{itemize}
  \item \textbf{``A banda está saturada, mais CO$_2$ não faz nada''}:
        o argumento do centro ignora as asas E a altura de emissão —
        T1 e N1 desmontam com números.
  \item \textbf{``1 K não é nada''}: 1\,K GLOBAL move distribuições
        inteiras — os extremos (Cap.~4: $+7\%$/K) escalam muito mais
        rápido que a média.
  \item \textbf{``Uma década fria refuta''}: detecção é multidecadal
        (Caso~4); ruído interno esconde o sinal por 30--40 anos.
  \item \textbf{``É o Sol''}: a estratosfera ESFRIA (T5) — o Sol
        aqueceria tudo junto; a assinatura vertical decide.
  \item \textbf{``Estabilizar emissões estabiliza o clima''}: TCRE
        exige líquido ZERO; estabilizar emissões só estabiliza a
        TAXA de aquecimento.
\end{itemize}

\subsection*{Demonstração de baixo custo}

Espectro de OLR medido por satélite (IRIS/AIRS) com o ``dente'' do
CO$_2$ — a estufa fotografada. Série do Niño-3.4 sobre a série global
de temperatura: o ruído do Caso~4 na vida real. Keeling completo
(dente de serra sazonal $+$ tendência): dois fenômenos do curso num
gráfico.

\subsection*{Problemas sugeridos do livro (módulo extra — sem seção correspondente no livro; usar os complementares)}

Aquecimento: $\Delta F$ para concentrações dadas; aquecimento de
equilíbrio por sensibilidade. Aprofundamento: orçamentos para metas;
expansão térmica. Síntese: ``liste três previsões do mecanismo estufa
feitas ANTES de serem observadas'' (estratosfera fria, noites
aquecendo mais, Ártico amplificado).
\emph{Confira a numeração na sua cópia (v1.02b).}

\subsection*{Uso do notebook \texttt{cap23\_mudanca.ipynb}}

\begin{center}
\begin{tabular}{lll}
\hline
Caso & Conteúdo & Momento sugerido \\
\hline
1 & forçante log e cenários por sensibilidade & Aula 1 \\
2 & duas caixas: aquecimento comprometido & Aula 1 \\
3 & TCRE e orçamento de carbono & Aula 2 \\
4 & emergência do sinal no ruído & Aula 2 \\
\hline
\end{tabular}
\end{center}

O Caso~2 congelando o CO$_2$ em 2024 rende a melhor discussão do
semestre: a dívida térmica como consequência direta do Cap.~3.

\subsection*{Exercícios complementares e soluções}

Nas notas: T1--T5 e N1--N4. Soluções em \texttt{cap23\_solucoes.ipynb}
(\emph{não distribuir}): T2 com \texttt{sympy} (escalas de tempo); N1
recupera o 5,35 de um modelo de banda; N3 compara trajetórias de
mesmo acumulado. Pares: T1$\leftrightarrow$N1, T2$\leftrightarrow$N2,
T4$\leftrightarrow$N3.

\subsection*{Conexões com o resto do curso}

TODAS, de propósito: estufa (Cap.~2), vapor e extremos (Cap.~4),
nuvens (Caps.~6--7), satélites que medem o dente do CO$_2$ (Cap.~8),
jato e bloqueios (Cap.~11), ciclones num mundo quente (Caps.~13, 16),
oceano lento (Cap.~3), modelos e ensembles (Cap.~20), clima natural
como régua (Cap.~21). O capítulo é a prova de que o curso forma
leitores qualificados do IPCC.


\end{document}
